\documentclass[11pt,letterpaper]{article}
\usepackage{cmap}
\usepackage[english]{babel}
\usepackage[margin=1in]{geometry}
\usepackage{amsmath,amssymb,amsthm}
\usepackage[hidelinks]{hyperref}
\hypersetup{
 pdftitle={Near-Balanced Permutations at Every Size},
 pdfauthor={Seonghyeon Bak; Jihoon Kim; Geonwoo Kim; Woojun Jeong},
 pdfsubject={Explicit permutation-pattern discrepancy and exact rectangular profiles},
 pdfkeywords={permutation patterns, balanced permutations, profile discrepancy, explicit construction}
}
\newtheorem{theorem}{Theorem}[section]
\newtheorem{definition}[theorem]{Definition}
\newtheorem{lemma}[theorem]{Lemma}
\newtheorem{proposition}[theorem]{Proposition}
\newtheorem{corollary}[theorem]{Corollary}
\newtheorem{remark}[theorem]{Remark}
\newcommand{\Ssym}{\mathfrak S}
\newcommand{\cB}{\mathcal B}
\title{Near-Balanced Permutations at Every Size}
\author{%
 Seonghyeon Bak\textsuperscript{1,\ensuremath{\dagger}}\quad
 Jihoon Kim\textsuperscript{2,\ensuremath{\dagger},*}\quad
 Geonwoo Kim\textsuperscript{3,\ensuremath{\ddagger}}\quad
 Woojun Jeong\textsuperscript{4,\ensuremath{\ddagger}}
}
\date{}
\begin{document}
\maketitle
\begingroup
\renewcommand{\thefootnote}{}
\begin{NoHyper}
\footnotetext{%
\textsuperscript{1}\,Department of Computer Science and Engineering,
Seoul National University, Seoul, Republic of Korea.\\
\textsuperscript{2}\,CAPP Lab, Department of Electrical and Computer Engineering,
Seoul National University, Seoul, Republic of Korea.\\
\textsuperscript{3}\,School of Undergraduate Studies,
Daegu Gyeongbuk Institute of Science and Technology, Daegu 42988, Republic of Korea.\\
\textsuperscript{4}\,Kyonggi University, Republic of Korea.\\
Seonghyeon Bak: \texttt{coshaman@snu.ac.kr};
Jihoon Kim: \texttt{jihoon@capp.snu.ac.kr}.\\
Geonwoo Kim: \texttt{koi@dgist.ac.kr};
Woojun Jeong: \texttt{jwj7140@gmail.com}.\\
\textsuperscript{\ensuremath{\dagger}}\,Seonghyeon Bak and Jihoon Kim contributed equally.\quad
\textsuperscript{\ensuremath{\ddagger}}\,Geonwoo Kim and Woojun Jeong contributed equally.\\
\textsuperscript{*}\,Corresponding author: Jihoon Kim.
}
\end{NoHyper}
\addtocounter{footnote}{-1}
\endgroup

\begin{abstract}
For $\pi\in\mathfrak S_N$, let $\delta_{\pi,k}$ be the largest deviation
of a length-$k$ pattern count from $\binom Nk/k!$, and let $\delta_k(N)$
be its minimum over $\mathfrak S_N$. Beniamini, Lavee, and Linial proved
$\delta_k(N)=\Omega_k(N^{k-1})$ for fixed $k\ge4$ and asked whether this
exponent is sharp for larger $k$. We construct an explicit family
$(\Pi_N)$, independent of $k$, with
$\delta_{\Pi_N,k}<14k^3N^{k-1}/(3k!)$ for all $2\le k\le N$.
Thus $\delta_k(N)=\Theta_k(N^{k-1})$ for every fixed $k\ge4$.
We also determine the sharp growing-pattern range of this family:
relative profile-density balance holds if and only if $k=o(N^{1/3})$.
At $k\sim cN^{1/3}$ the increasing pattern has relative density
$e^{-c^3/6}$. The construction completes a signed two-sheet rectangle
by a mechanical partial row; the proofs isolate connected coordinate
collisions after cancelling the contributions of isolated pairs.
Exact rectangular profiles also identify the leading discrepancy removable
by bounded-prefix rows and bounded insertions. For each fixed $k\ge4$
and aspect $\theta>0$, a nonconstructive realization theorem attains the
standard-sheet LP distance $q_k(\theta)/[2k!(k-2)!]$. Certificates determine
optimal aspects for $k=4,\ldots,8$, including a nonsquare optimum at $k=6$.
A separate explicit family with periodic row orientations and $12$ or $17$
insertions attains four-pattern constant $1/192$, matching a lower bound
within that completion class, not over arbitrary permutations.
\end{abstract}

\noindent\textbf{Keywords.} permutation patterns; balanced permutations;
profile discrepancy; explicit construction; Erd\H{o}s--Szekeres permutation.

\noindent\textbf{2020 Mathematics Subject Classification.}
Primary 05A05; Secondary 05D99.

\section{Introduction}

For pairwise distinct real numbers $z_1,\ldots,z_k$, let
$\operatorname{st}(z_1,\ldots,z_k)\in\Ssym_k$ be the permutation that
records their relative order.  If $\pi\in\Ssym_N$ and
$i_1<\cdots<i_k$, the induced pattern is
$\operatorname{st}(\pi(i_1),\ldots,\pi(i_k))$.  Write $\#_\tau(\pi)$
for the number of $k$-subsets of positions inducing $\tau\in\Ssym_k$.
For example, $\operatorname{st}(4,1,3)=312$.
The $k$-profile of $\pi$ is the vector of these $k!$ counts.  We study
\[
 \delta_{\pi,k}=
 \max_{\tau\in\Ssym_k}\left|\#_\tau(\pi)-\frac1{k!}\binom Nk\right|,
 \qquad
 \delta_k(N)=\min_{\pi\in\Ssym_N}\delta_{\pi,k}.
\]
Thus $\pi$ is $k$-balanced precisely when $\delta_{\pi,k}=0$.

A uniformly random permutation has expected profile $\binom Nk/k!$, while
quasirandom permutation sequences approach these densities asymptotically
\cite{Cooper,KralPikhurko}. Related frameworks include permutation limits
\cite{HoppenEtAl} and asymptotic profile symmetry \cite{CooperPetrarca}.
Here the objective is a finite, simultaneous $\ell_\infty$
approximation to the uniform profile.  Joy Cooper and Dukes
\cite[equation~(1.2)]{CooperDukesFuzzy} define the fuzzy lift
$P_\sigma^{\uparrow n}$ as a normalized sum of order-preserving embedding
matrices.  They study rational combinations of these lifts that are constant
matrices.  The same binomial weights appear in our insertion expansion,
but those formal identities do not provide the every-size discrepancy bound
proved here: we realize centred corrections by unweighted points in one
permutation.
Exact balance is exceptionally rigid.  Beniamini, Lavee, and Linial
proved that no permutation is $4$-balanced and obtained the lower bound
$\delta_k(N)=\Omega_k(N^{k-1})$ for every fixed $k\ge4$
\cite[Theorems~6 and 7]{BLL}.  Their two-sided Erd\H{o}s--Szekeres
construction matches this order for $k=4$ along orders $2m^2$, but the
sharp exponent for each fixed $k>4$ was left open.
An explicit all-size baseline follows from Cooper's estimate
\cite[Theorem~3.1 and the proof of Proposition~3.4]{Cooper}:
for every fixed $k\ge2$, every $N>k$, and every $\pi\in\Ssym_N$,
\[
 \delta_{\pi,k}\le C_kD(\pi)N^{k-1},\qquad
 D(\pi):=\max_{\substack{I,J\subseteq[N]\\ I,J\text{ intervals}}}
 \left|\#\bigl(\pi(I)\cap J\bigr)-\frac{|I||J|}{N}\right|,
\]
Here intervals are nonwrapping; $D_{\rm cyc}$ denotes the same maximum
over cyclic intervals in $\mathbb Z_N$.  Splitting each cyclic interval into
at most two nonwrapping intervals gives $D_{\rm cyc}\le4D$.
Cooper's decoupling estimate is $O_k((D_{\rm cyc}+1)N^{k-1})$.
Taking $I=[N-1]$ and $J=\{\pi(N)\}$ gives $D\ge(N-1)/N$,
which absorbs the additive $1$ and yields the displayed finite bound.  The S\'os permutation
$\sigma_N^{(\alpha)}$ ranking the fractional parts $\{\alpha\},\ldots,\{N\alpha\}$, with
$\alpha=(\sqrt5-1)/2$, has $D(\sigma_N^{(\alpha)})=O_\alpha(\log N)$
\cite[Theorem~3.24]{CooperArithmetic}.  Hence
$\delta_{\sigma_N^{(\alpha)},k}=O_{k,\alpha}(N^{k-1}\log N)$ at every size.
The issue is to remove this logarithm while retaining all patterns and
all permutation sizes.

Throughout, $N$ denotes permutation size and $k$ denotes pattern length.

\begin{theorem}\label{thm:main}
There is an explicit family $(\Pi_N)_{N\ge1}$, independent of $k$, with the
following properties.
\textup{(i)} For all $2\le k\le N$,
\[
 \delta_{\Pi_N,k}<\frac{14k^3}{3k!}N^{k-1}.
\]
Consequently $\delta_k(N)=\Theta_k(N^{k-1})$ for every fixed $k\ge4$.
\par\noindent\textup{(ii)} For $2\le k=k(N)\le N$, relative profile-density balance
\[
 \max_{\tau\in\Ssym_k}
 \left|\frac{k!\,\#_\tau(\Pi_N)}{\binom Nk}-1\right|\longrightarrow0
\]
holds if and only if $k=o(N^{1/3})$. If $k/N^{1/3}\to c\in[0,\infty)$,
the increasing pattern has relative density tending to $e^{-c^3/6}$;
if $k/N^{1/3}\to\infty$, its relative density tends to zero.
\end{theorem}

Theorem~\ref{thm:main} closes the fixed-$k$ exponent question of
Beniamini--Lavee--Linial and removes the logarithm from Cooper's all-size
comparison. Part (ii) determines exactly when this family approximates
every shrinking target $1/k!$ relatively.
Its threshold is specific to $\Pi_N$, not a bound for all permutation families.

The complete two-sheet rectangle is inherited from \cite{BLL}.  The
obstacle to an every-size construction is its boundary: filling a gap of
$O(\sqrt N)$ points can affect $\Theta_k(N^{k-1/2})$ subsets, more than the
allowed discrepancy.  We use a partial row whose columns form a mechanical
word.  Sign reversal makes isolated coordinate-tie resolutions unbiased;
a coupling to independent tie-breaking leaves only interacting collisions.
Prefix balance then controls the centred profile of the large clean boundary
class, rather than bounding its cardinality.  One northeast point handles
odd sizes. This proves part (i) in
Sections~\ref{sec:construction}--\ref{sec:odd-orders}.

The questions of pattern growth and removable leading error are linked by
the exact rectangle formula in Section~\ref{sec:rectangles}. A connected
collision expansion cancels arbitrary collections of isolated pairs and
bounds the remaining relative error by $O(k^3/N)$ on complete squares.
An exact increasing-pattern count gives the matching obstruction at
$k\asymp N^{1/3}$. The same formula identifies three local statistics
that describe the fixed-$k$ leading vector.

Section~\ref{sec:sharp-four} turns that vector into a completion optimum.
A bounded-prefix row contributes a term depending only on the final pattern
value, and boundedly many insertions contribute sums over position--value
pairs.  We prove that every bounded centred correction of this additive
form can be realized.  The remaining discrepancy is therefore exactly the
$\ell_\infty$ distance $q_k(\theta)$ of the core vector from the additive
subspace, with normalization $2k!(k-2)!$.  Exact certificates show that square
aspect is optimal in this class for $k=4,5,7,8$, but not for $k=6$: its
minimizing aspects are $2-2\sqrt3/3$ and its reciprocal.  Its family $\Xi_N^{(k,\theta)}$ is nonconstructive and may depend on both
fixed parameters, unlike $\Pi_N$.

The completion class matters even for $k=4$.  Its standard-orientation
optimum is $1/144$, but periodically reversing rows adds a local bond term to
the core vector.  We exploit it to construct a separate explicit family
$(\Sigma_N)$ with $12$ or $17$ insertions and discrepancy
$N^3/192+o(N^3)$.  A ten-pattern functional annihilates centred additive
corrections and proves optimality over all fixed-period orientations and
positive limiting aspects in the stated completion framework.  No constant
optimum here is asserted over arbitrary permutations.
\section{The every-size construction}\label{sec:construction}

Write $[s]=\{1,\ldots,s\}$ and compare tuple-valued keys
lexicographically.  For even size, write $N=2E$ and set
\[
 a=\lfloor\sqrt E\rfloor,
 \qquad b=\left\lfloor\frac Ea\right\rfloor,
 \qquad r=E-ab.
\]
Then
\begin{equation}\label{eq:shape}
 0\leq r<a,
 \qquad a\leq b\leq a+2.
\end{equation}
If $r>0$, define
\[
 R=R_{b,r}
   =\left\{\left\lceil\frac{jb}{r}\right\rceil:1\leq j\leq r\right\}
   \subseteq[b];
\]
When $r>0$, \eqref{eq:shape} gives $b/r>1$, so the $r$ ceiling values
are distinct and $|R|=r$.  Put $R=\varnothing$ when $r=0$.  Our cell set is
\[
 F_E=([a]\times[b])\,\cup\,(\{a+1\}\times R),
\]
which has exactly $E$ cells.  The first rectangle is the \emph{core};
the cells in the last row are \emph{partial}.
For example, $E=7$ gives $(a,b,r)=(2,3,1)$, $R=\{3\}$, and hence a
$2\times3$ core together with the single partial cell $(3,3)$.

Place two coloured points above each cell:
\[
 Q_E=F_E\times\{-1,+1\}.
\]
For $p=(x,y,c)$, order these $2E$ points by the two injective keys
\begin{align}
 K_X(p)&=(x,-cy,c),
 &K_Y(p)&=(y,cx,c).                         \label{eq:keys}
\end{align}
Listing the points in $K_X$-order and recording their ranks in
$K_Y$-order defines $\Pi_{2E}\in\Ssym_{2E}$.  More generally, for any
finite cell set $F\subset\mathbb Z^2$, write $\Pi(F)$ for the permutation
of $F\times\{-1,+1\}$ induced by these same two key orders.  The final
coordinate keeps both keys injective even at zero or negative coordinates;
the rotation interpretation below concerns the positive-coordinate set $F_E$.
For the preceding $E=7$ example this gives
\[
 \Pi_{14}=(12,7,3,2,6,11,13,8,4,1,5,10,14,9).
\]

On the positive-coordinate set $F_E$, the final coordinate is redundant.
Here the two orders also have an exact geometric realization.  Put $M=\max\{a+1,b\}$ and choose
$0<\varepsilon<\arctan(1/(2M))$.  Rotate the copy of
colour $c$ through angle $c\varepsilon$; its coordinates become
\[
 (x\cos\varepsilon-cy\sin\varepsilon,\quad
  y\cos\varepsilon+cx\sin\varepsilon).
\]
A unit difference between original $x$-coordinates contributes
$\cos\varepsilon$, while the perturbation difference has magnitude at
most $2b\sin\varepsilon$; hence unequal $x$-coordinates keep their
order.  The analogous bound $2(a+1)\sin\varepsilon$ handles unequal
$y$-coordinates.  Original ties are resolved by $-cy$ and $cx$.
Therefore the rotated coordinate orders are precisely
\eqref{eq:keys}.  We write
$\mathrm{ES}^{\pm}(a,b):=\Pi([a]\times[b])$ for this complete-core
permutation.  It is Beniamini--Lavee--Linial's sufficiently-small-angle order
type; if $r=0$, it is the whole construction \cite[Definition~5.6]{BLL}.

\begin{figure}[t]
\centering
\setlength{\unitlength}{1pt}
\begin{picture}(440,128)
  \put(0,116){\makebox(170,8){\small (a) two sheets over a partial rectangle}}
  \put(18,16){\line(1,0){62}}
  \put(18,16){\line(0,1){76}}
  \put(80,16){\line(0,1){76}}
  \put(18,92){\line(1,0){62}}
  \put(34,28){\circle*{3}} \put(40,22){\circle{4}}
  \put(64,28){\circle*{3}} \put(70,22){\circle{4}}
  \put(34,53){\circle*{3}} \put(40,47){\circle{4}}
  \put(64,53){\circle*{3}} \put(70,47){\circle{4}}
  \put(34,78){\circle*{3}} \put(40,72){\circle{4}}
  \put(64,78){\circle*{3}} \put(70,72){\circle{4}}
  \put(94,78){\circle*{3}} \put(100,72){\circle{4}}
  \multiput(97,61)(0,5){8}{\line(0,1){2}}
  \put(30,4){\makebox(40,8){\scriptsize $[a]\times[b]$ core}}
  \put(84,50){\makebox(36,8){\scriptsize partial row}}
  \put(5,94){\makebox(18,8){\scriptsize $y$}}
  \put(82,8){\makebox(18,8){\scriptsize $x$}}
  \put(18,104){\circle*{3}} \put(23,100){\makebox(38,8)[l]{\scriptsize $c=+1$}}
  \put(84,104){\circle{4}} \put(89,100){\makebox(38,8)[l]{\scriptsize $c=-1$}}

  \put(172,116){\makebox(268,8){\small (b) cancellation hierarchy}}
  \put(190,92){\makebox(70,8){\scriptsize collision-free}}
  \put(207,66){\circle*{3}} \put(219,78){\circle{4}}
  \put(232,58){\circle*{3}} \put(247,88){\circle{4}}
  \put(190,40){\makebox(70,8){\scriptsize pattern-uniform}}

  \put(274,92){\makebox(78,8){\scriptsize isolated $x$-pair}}
  \put(298,61){\circle*{3}} \put(298,82){\circle{4}}
  \put(298,63){\line(0,1){17}}
  \put(318,55){\circle*{3}} \put(337,88){\circle{4}}
  \put(270,40){\makebox(88,8){\scriptsize uniform after colours}}

  \put(366,92){\makebox(66,8){\scriptsize interacting wedge}}
  \put(386,70){\circle*{4}} \put(386,88){\circle{4}}
  \put(408,70){\circle*{3}}
  \put(386,72){\line(0,1){14}} \put(388,70){\line(1,0){18}}
  \put(362,40){\makebox(76,8){\scriptsize discrepancy source}}
  \put(186,19){\makebox(245,10){\scriptsize coordinate equalities are shown by line segments}}
\end{picture}
% Alt text: a two-by-three core plus one partial cell, with two signed
% points over each cell, followed by collision-free, isolated-pair, and
% interacting-wedge configurations illustrating the cancellation hierarchy.
\caption{Construction and cancellation mechanism.  (a) The case $E=7$:
a $2\times3$ core and one partial cell, with a positive (filled) and
negative (open) point over each cell.  The keys in \eqref{eq:keys} are the
coordinate orders after opposite small rotations of the two sheets.
(b) Collision-free classes and isolated same-axis pairs are exactly
pattern-uniform after summing colours (Lemma~\ref{lem:core-uniform}); only
interacting equalities, such as the displayed wedge, enter the discrepancy
bound.}
\label{fig:construction}
\end{figure}

For odd $N>1$, obtain $\Pi_N$ from $\Pi_{N-1}$ by appending a point that
is both rightmost and highest.  Let $\Pi_1$ be the unique permutation of
size one.

\section{A balanced partial row}

The set $R$ distributes its $r$ columns evenly through $[b]$.

\begin{lemma}[prefix balance]\label{lem:prefix}
For every integer $0\leq z\leq b$,
\[
 |R\cap[z]|=\left\lfloor\frac{rz}{b}\right\rfloor.
\]
In particular, if $d_y=\mathbf1_R(y)-r/b$, then
\[
 \left|\sum_{y=1}^z d_y\right|<1
 \quad (0\leq z<b),
 \qquad
 \sum_{y=1}^b d_y=0.
\]
\end{lemma}

\begin{proof}
For $r>0$, the inequality $\lceil jb/r\rceil\leq z$ is equivalent to
$j\leq rz/b$.  Counting the eligible integers $j$ gives the formula.
The case $r=0$ is immediate.
\end{proof}

\begin{lemma}[integrated mechanical discrepancy]\label{lem:mechanical-earthmover}
For $0\leq r<b$, the mechanical row satisfies
\[
 W(R):=\sum_{z=1}^{b-1}
 \left||R\cap[z]|-\frac{rz}{b}\right|
 =\frac{b-\gcd(b,r)}2\leq\frac{b-1}2.
\]
\end{lemma}

\begin{proof}
Lemma~\ref{lem:prefix} makes the $z$th summand the fractional part
$\{rz/b\}$.  If $g=\gcd(b,r)$, the nonzero residues of $rz$ modulo $b$
consist of each of $g,2g,\ldots,b-g$ repeated $g$ times.  Their normalized
sum is $(b-g)/2$.
\end{proof}

For $1\leq t\leq k$ and $1\leq y\leq b$, use the convention that an
out-of-range binomial coefficient is zero and put
\[
 f_t(y)=\binom{y-1}{t-1}\binom{b-y}{k-t}.
\]
This counts the $(k-1)$-subsets of $[b]\setminus\{y\}$ in which $y$ has
rank $t$ after it is inserted.

\begin{lemma}[rank balance]\label{lem:rank-balance}
For the mechanical row $R$,
\begin{equation}\label{eq:rank-balance}
 \left|\sum_{y\in R} f_t(y)-\frac rb\binom bk\right|
 \leq\binom{b-1}{k-1}.
\end{equation}
\end{lemma}
\begin{proof}
Vandermonde gives $\sum_y f_t(y)=\binom bk$.
Put $g_y={\bf1}_R(y)-r/b$ and $G_z=\sum_{y\le z}g_y$.
Lemma~\ref{lem:prefix} gives $-1<G_z\le0$ and $G_b=0$.
The nonnegative weight $f_t$ is unimodal: where adjacent terms are positive,
$f_t(y+1)/f_t(y)=y(b-y-k+t)/((y-t+1)(b-y))$ is at least one exactly when
$y\le(t-1)b/(k-1)$. Also $f_t(y)\le\binom{b-1}{k-1}$, since it counts
only the $(k-1)$-subsets avoiding $y$ that put $y$ at rank $t$.
In $\sum_y g_yf_t(y)=-\sum_{z<b}G_z(f_t(z+1)-f_t(z))$,
the total increase and total decrease are each bounded by this
maximum, and $0\le-G_z<1$. This proves the claim, including $r=0$.
\end{proof}

\section{Exact cancellation in the core}

Consider first the coloured complete rectangle
$[a]\times[b]\times\{-1,+1\}$ with the orders
\eqref{eq:keys}.  For a $k$-subset $S$, let $r_x(S)$ and $r_y(S)$ be
the numbers of distinct $x$- and $y$-coordinates, and define its
\emph{collision defect}
\[
 d(S)=(k-r_x(S))+(k-r_y(S)).
\]
The three configurations in Figure~\ref{fig:construction}(b) explain the
threshold.  If all coordinates are distinct, every one of the $2^k$ colour
words gives the same two strict coordinate orders.  If exactly one
$x$-coordinate is repeated while all ordinates are distinct, then for any
prescribed pattern the coordinate assignment is forced; exactly two of the
four colour pairs resolve the tied pair in the required order, and the other
$k-2$ colours are free, giving $2^{k-1}$ realizations of every pattern.  If an
$x$-tie and a $y$-tie share a point, its single colour must satisfy two local
resolutions, which may agree or conflict.  These interacting equalities are
the first classes that need an estimate.

\begin{lemma}[same-axis pair-collision cancellation]
\label{lem:core-uniform}
In the complete rectangle, fix an $x$-coordinate multiset and a set of
distinct $y$-coordinates.  Every pattern occurs equally often among the
corresponding coloured subsets if every $x$-multiplicity is at most two.
Symmetrically, the same holds
for a fixed set of distinct $x$-coordinates and a fixed $y$-coordinate
multiset whose multiplicities are at most two.  In particular, the
layers $d=0$ and $d=1$ separately have equal counts for all patterns.
\end{lemma}

\begin{proof}
Suppose the $x$-multiset has $q$ repeated pairs and all ordinates are
distinct.  The pairs occupy disjoint adjacent $X$-blocks.  A prescribed
pattern fixes the unordered ordinate pair in each block and the ordinate
of each singleton, hence the uncoloured cell set.  Reversing both colours
negates the two distinct tie keys $(-cy,c)$ and reverses their comparison,
so either resolution has exactly two colour choices.  The $q$ pairs use
disjoint colour variables and the $k-2q$ singleton colours are free.
Thus exactly $2^q2^{k-2q}=2^{k-q}$ colourings realize each pattern.
Exchange axes and use $(cx,c)$ for the other assertion.  Summing over
coordinate choices proves uniformity; $q=0$ and $q=1$ give the stated
defect layers.
\end{proof}

The following transfer statement does not require a rectangle or positive
coordinates.  For a finite cell set $F\subset\mathbb Z^2$, put $E=|F|$.  Write $r_x$ and $c_y$ for its
row and column degrees, and set
\[
 R_3=\sum_xr_x^3,\qquad C_3=\sum_yc_y^3,
 \qquad M_{11}=\sum_{(x,y)\in F}r_xc_y.
\]
For a vector indexed by $\Ssym_k$, let $\|\cdot\|_{\rm cen}$ denote
maximum distance from its coordinate average.

\begin{lemma}[independent-jitter transfer]\label{lem:jitter}
Let $P_F$ be the $k$-profile vector of $\Pi(F)$.  For each ordered
cell tuple $z\in F^k$, independently refine every equality class of its
$x$- and $y$-coordinates by a uniform label order, and let
$q_z(\tau)$ be the probability that the two refined orders induce
$\tau$.  Define the unnormalized jitter-mass vector by
\[
 J_F(\tau)=2^k\sum_{z\in F^k}q_z(\tau).
\]
Then
\[
 \|k!P_F-J_F\|_{\rm cen}\leq 2^k I_k(F),
 \tag{4.1}\label{eq:jitter-transfer}
\]
where, for $k\geq3$,
\[
 I_k(F)=\binom{k}{3}(R_3+C_3)E^{k-3}
 +k(k-1)(k-2)M_{11}E^{k-3}
 +\binom{k}{2}E^{k-1}.
 \tag{4.2}\label{eq:interaction-moment}
\]
For $k=2$, the right side of \eqref{eq:jitter-transfer} is interpreted
with $I_2(F)=E$.
\end{lemma}

\begin{proof}
For $z=(z_1,\ldots,z_k)\in F^k$ and
$c=(c_1,\ldots,c_k)\in\{-1,+1\}^k$, let
$\operatorname{pat}(z,c)$ be the pattern induced by the
$K_X$- and $K_Y$-orders of the labelled signed points
$p_i=(z_i,c_i)$ when they are pairwise distinct.
\[
 A_z(\tau)=\sum_{c\in\{-1,+1\}^k}
 {\bf1}\{p_1,\ldots,p_k\text{ are distinct and }
             \operatorname{pat}(z,c)=\tau\},
 \qquad B_z(\tau)=2^kq_z(\tau).
\]
Every coloured $k$-subset is represented by its $k!$ label orders, while
$q_z$ is a probability vector.  Thus, componentwise,
\[
 k!P_F=\sum_{z\in F^k}A_z,
 \qquad J_F=\sum_{z\in F^k}B_z.
\]

Join two labels when their cells have the same $x$-coordinate or the same
$y$-coordinate.  Call $z$ good when this union graph is a matching and no
two labels occupy the same cell.  If the graph has $m$ edges, every choice
of its $m$ independent jitter resolutions has probability $2^{-m}$.  We
spell out the corresponding colour count.  On an $x$-edge $\{i,j\}$ the
ordinates are distinct, so the relative $Y$-order is already fixed.
The tie keys are $u_i=(-c_i y_i,c_i)$ and $u_j=(-c_j y_j,c_j)$.
They are distinct: equality would force $c_i=c_j$ and then $y_i=y_j$.
The fixed-point-free involution $(c_i,c_j)\mapsto(-c_i,-c_j)$ sends
$(u_i,u_j)$ to $(-u_i,-u_j)$, reversing their strict lexicographic
comparison.  It therefore pairs each resolution with its opposite, so
exactly two of the four colour pairs realize either order.  This argument
includes negative and zero ordinates and ties in the first key component.
On a $y$-edge, the distinct abscissae give the identical argument for
$(c_i x_i,c_i)$ and $(c_j x_j,c_j)$.  Because the union graph is a matching, these
edge constraints use disjoint colour variables; because repeated cells are
excluded, no edge is simultaneously an $x$- and a $y$-edge.  The $k-2m$ singleton colours are free, so each joint resolution has
$2^m2^{k-2m}=2^{k-m}$ realizing words.  The $2^m$ resolutions partition
all $2^k$ colour words.  Each has jitter probability $2^{-m}$ and hence
mass $2^{k-m}$ in $B_z=2^kq_z$ as well.  Thus $A_z=B_z$ for every good $z$.

Disagreement can therefore occur only when an axis contains three chosen
labels, one label is incident to an $x$-edge and a $y$-edge with different
neighbours, or a label pair occupies the same cell.  For the first event,
choose the three responsible labels and the common row (or column), then
choose their cells and the remaining cells; this gives
$\binom{k}{3}(R_3+C_3)E^{k-3}$.  For a mixed wedge, choose its ordered
centre and two neighbours and then its centre cell; the two incident
choices give $r_xc_y$, hence
$k(k-1)(k-2)M_{11}E^{k-3}$.  A repeated cell is bounded by choosing its
label pair, the cell, and the remaining cells, giving
$\binom{k}{2}E^{k-1}$.  These are union bounds, so overlaps need not be
subtracted.  For $k=2$, only the repeated-cell event remains.

Finally, if nonnegative vectors $u,v\in\mathbb R^n$ have coordinate sums
$U,V\le M$, then, coordinate by coordinate,
\[
 \left|(u_i-v_i)-\frac{U-V}{n}\right|
 \le \max\left\{U\left(1-\frac1n\right)+\frac Vn,
                  V\left(1-\frac1n\right)+\frac Un\right\}
 \le M.
\]
For every bad tuple, $A_z$ and $B_z$ are nonnegative vectors with
coordinate sums at most $2^k$ (the sum for $B_z$ is exactly $2^k$).
Taking $n=k!$ and $M=2^k$, then summing over the bad-tuple union bound,
proves \eqref{eq:jitter-transfer}.
\end{proof}

For the complete $a\times b$ rectangle, summing over all ordered cell tuples
is equivalent to choosing the labelled $x$-coordinates independently and
uniformly from $[a]$, and the labelled $y$-coordinates independently and
uniformly from $[b]$.  Randomly refining ties therefore makes each global
order uniform; independence of the coordinate choices and refinements makes
the two orders independent.  Hence $J_F$ is exactly uniform.  With $A=ab$
and $\eta=b/a$, Lemma~\ref{lem:jitter} therefore gives
\begin{equation}
 \delta_{\Pi([a]\times[b]),k}
 \leq \frac{2^k}{k!}A^{k-1}K_{k,\eta},
 \qquad
 K_{k,\eta}=\binom{k}{3}(\eta+\eta^{-1})
             +\binom{k}{2}(2k-3).
 \tag{4.3}\label{eq:jitter-rectangle}
\end{equation}

\section{Balanced boundary completion}\label{sec:boundary}

We now count subsets that use the partial row.  A subset with at least
two partial-row points belongs to a family of size at most
\begin{equation}\label{eq:two-partial}
 \binom{2r}{2}\binom{2E}{k-2}
 \leq\frac{2^{k-1}}{(k-2)!}E^{k-1}.
\end{equation}

Consider a subset with exactly one partial-row point.  In this count a
\emph{coordinate collision} means equality of original $x$- or $y$-coordinates
among selected signed points.  The two opposite-colour points over one cell are
distinct labelled points, but selecting both creates a collision on each axis;
selecting the same labelled point twice is never allowed.  For $k\geq3$, a
labelled $(k-1)$-tuple with a repeated coordinate realizes one of
$J_k=(k-1)(k-2)$ elementary equality events.  For a fixed event there
are at most $bA^{k-2}$ core cell assignments, since $a\leq b$.
Dividing by $(k-1)!$, choosing colours and one of the $2r$ partial
points, and using $rb\leq ab=A$, gives
\[
 \frac{2^kJ_k}{(k-1)!}A^{k-1}.
\]
This single event count includes repetitions on either axis.  For $k=2$
this repeated-coordinate subclass is empty.
There is one additional collision possibility: a core point may share
the partial point's $y$-coordinate.  A union bound gives at most
\begin{equation}\label{eq:shared-y}
 (2r)(2a)\binom{2ab}{k-2}
 \leq\frac{2^k}{(k-2)!}E^{k-1}
\end{equation}
subsets of this kind.  The partial row has a new $x$-coordinate, so
there is no analogous shared-$x$ case.

It remains to analyze subsets with exactly one partial-row point and no
coordinate collision.  Their exact number is
\[
 2^k r (k-1)!\binom a{k-1}\binom{b-1}{k-1}.
\]
Indeed, choose the signed partial point, the distinct core abscissae, the
core ordinates other than its ordinate, their bijection, and the $2^{k-1}$
core colours.  Thus a crude cardinality bound is insufficient: when
$r\asymp a\asymp b\asymp E^{1/2}$, this class has
$\Theta_k(E^{k-1/2})$ members.  The next lemma uses the bounded rank
discrepancy of the mechanical row to show that its \emph{centred profile},
rather than its size, is only $O_k(E^{k-1})$.

For $k\geq2$, define $m_k=1$ when $k=2$ and $m_k=2$ when $k\geq3$.

\begin{lemma}[clean boundary profile]\label{lem:boundary-profile}
Let $A_\tau$ be the number of collision-free $k$-subsets that contain
exactly one partial-row point and induce $\tau$.  If
$t=\tau(k)$, then
\begin{equation}\label{eq:A-tau}
 A_\tau
 =2^k\binom{a}{k-1}
   \sum_{y\in R}
     \binom{y-1}{t-1}\binom{b-y}{k-t}.
\end{equation}
Consequently,
\begin{equation}\label{eq:boundary-explicit}
 \left|A_\tau-\frac1{k!}\sum_{\sigma\in\Ssym_k}A_\sigma\right|
 \leq2^km_k\binom{a}{k-1}\binom{b-1}{k-1}
 \leq\frac{2^km_k}{((k-1)!)^2}E^{k-1}.
\end{equation}
\end{lemma}

\begin{proof}
For $k=2$ the argument below has one core point and gives the same formula
directly, with $m_2=1$.  We therefore treat all $k\geq2$ together.
The partial row is last in $X$-order.  Fix its ordinate $y\in R$.
For it to have rank $t$ among the selected ordinates, choose $t-1$
ordinates below $y$ and $k-t$ above it.  Choose the other $k-1$
distinct abscissae from the $a$ core rows.  Once these coordinate sets
and a pattern $\tau$ with $\tau(k)=t$ are fixed, $\tau$ specifies a
unique bijection from the core abscissae to the non-special ordinates.
There are no ties, so all $2^k$ colourings induce the same pattern.
This proves \eqref{eq:A-tau}.

Write the sum in Lemma~\ref{lem:rank-balance} as
$r\binom bk/b+e_t$, where
$|e_t|\leq\binom{b-1}{k-1}\leq m_k\binom{b-1}{k-1}$.
Moreover, $\sum_t f_t(y)=\binom{b-1}{k-1}$, and Vandermonde's identity
also gives $k\binom bk/b=\binom{b-1}{k-1}$.  Hence
$\sum_{t=1}^k e_t=0$.  Each final rank occurs for exactly $(k-1)!$
patterns, so the class average has zero error term.  Multiplying the
bound on $e_t$ by $2^k\binom{a}{k-1}$ proves
\eqref{eq:boundary-explicit}.
\end{proof}

Put
\[
 \beta_k=1-\frac1{k!},\qquad J_k=(k-1)(k-2),
\]
and, for $\rho\geq1$, define
\begin{equation}\label{eq:C-rho}
 \begin{aligned}
 C_\rho(k)={}&\frac{2m_k}{((k-1)!)^2}
 +\frac{2K_{k,\rho}}{k!}\\
 &+\beta_k\left(
   \frac{2J_k}{(k-1)!}
  +\frac3{(k-2)!}\right)
  +\frac{\beta_k}{(k-1)!},\\
 K_{k,\rho}={}&\binom{k}{3}(\rho+\rho^{-1})
                 +\binom{k}{2}(2k-3).
 \end{aligned}
\end{equation}
The last term is reserved for odd padding.
The binomial convention makes every displayed class
count valid even when there are too few rows or columns; in particular, the
complete-core class is simply empty when $k>2ab$.  \begin{proof}[Proof of Theorem~\ref{thm:main} for even $N$]
If $k>2ab$, the class using only the $2ab$ core points is empty; throughout
this proof its centred contribution is then interpreted as zero.  Otherwise
the complete core contributes discrepancy at most
\eqref{eq:jitter-rectangle}.  The clean one-partial-point class
contributes at most \eqref{eq:boundary-explicit}.  Put every remaining
partial-row subset into $\cB$.  Equations
\eqref{eq:two-partial}--\eqref{eq:shared-y} give, with
$\rho_E=b/a$,
\[
 |\cB|\leq2^{k-1}\left(
  \frac{2J_k}{(k-1)!}
  +\frac3{(k-2)!}\right)E^{k-1}.
\]
Since the centered contribution of an arbitrary class of size $L$ is at
most $\beta_kL$, and $N=2E$, the three contributions give
\[
 \overline C_E(k)=C_{\rho_E}(k)
 -{\bf1}_{\{k>2ab\}}\frac{2K_{k,\rho_E}}{k!}
\]
and
\begin{equation}\label{eq:refined-even}
 \delta_{\Pi_N,k}
 \leq\left(\overline C_E(k)-\frac{\beta_k}{(k-1)!}\right)N^{k-1}.
\end{equation}
To normalize \eqref{eq:C-rho}, use $\beta_k\le1$ and set
$\eta=\rho+\rho^{-1}$.  The clean term is at most $2/k^2$ after
multiplication by $k!/k^3$; the bad and reserved terms sum to
$(2J_k+3(k-1)+1)/k^2=2-3/k+2/k^2$.
If $k\le2ab$, then $b-a\le2$ gives
$\eta-2=(b-a)^2/(ab)\le8/k$, and therefore
\[
\begin{aligned}
 \frac{k!}{k^3}C_{\rho_E}(k)
 &\le\frac{\eta(k-1)(k-2)}{3k^2}
 +\frac{(k-1)(2k-3)}{k^2}+2-\frac3k+\frac4{k^2}\\
 &\le\frac{14k^3-22k^2+k+16}{3k^3}<\frac{14}{3}.
\end{aligned}
\]
If $k>2ab$, omit the empty core term.  The remaining normalized sum is
at most $(2k^2-3k+4)/k^2<14/3$.  Thus
$\overline C_E(k)<14k^3/(3k!)$ in both cases.
\end{proof}

\section{Odd orders and completion of the proof}\label{sec:odd-orders}

\begin{lemma}[one-point padding]\label{lem:padding}
Let $2\leq k\leq M+1$.  Suppose $\pi\in\Ssym_M$ and
$\pi^+\in\Ssym_{M+1}$ is obtained by appending a
rightmost-and-highest point.  Write
\[
 D_k(\tau)=\#_\tau(\pi)-\frac1{k!}\binom Mk,
 \qquad
 D_k^+(\tau)=\#_\tau(\pi^+)-\frac1{k!}\binom{M+1}k,
 \qquad
 c_{M,k}=\frac1{k!}\binom M{k-1}.
\]
Set $D_{M+1}(\tau)=0$, since $\pi$ has no $(M+1)$-subsets, and
$D_1(1)=0$.
If $\tau(k)=k$, let $\tau'$ be obtained by deleting its last entry.  Then
\begin{equation}\label{eq:odd-recurrence}
 D_k^+(\tau)=D_k(\tau)+
 \begin{cases}
  D_{k-1}(\tau')+(k-1)c_{M,k},&\tau(k)=k,\\
  -c_{M,k},&\tau(k)\ne k.
 \end{cases}
\end{equation}
With the conventions $\delta_{\pi,1}=\delta_{\pi,M+1}=0$,
\begin{align}
 \delta_{\pi^+,k}
 &\leq \delta_{\pi,k}
   +\beta_k\binom{M}{k-1},
 \label{eq:padding-coarse}\\
 \delta_{\pi^+,k}
 &\leq \delta_{\pi,k}+\delta_{\pi,k-1}
   +\frac{k-1}{k!}\binom{M}{k-1}.
 \label{eq:padding-sharp}
\end{align}
\end{lemma}

\begin{proof}
The number $X_\tau$ of new copies is $\#_{\tau'}(\pi)$ when
$\tau(k)=k$, and zero otherwise.  Pascal's identity gives
\eqref{eq:odd-recurrence}.  Centering the range
$0\le X_\tau\le\binom M{k-1}$ at $c_{M,k}$ gives
\eqref{eq:padding-coarse}; taking absolute values in the recurrence gives
\eqref{eq:padding-sharp}.
\end{proof}

For odd $N\geq k+1$, let $M=N-1$.  Combining
\eqref{eq:padding-coarse} with the reserved term in the even estimate gives
the explicit inequality
\[
 \delta_{\Pi_N,k}
 \leq\left(\overline C_E(k)-\frac{\beta_k}{(k-1)!}\right)M^{k-1}
      +\beta_k\binom{M}{k-1}
 \leq \overline C_E(k)M^{k-1}
 \leq \overline C_E(k)N^{k-1}.
\]
If $N=k$ is odd, then
$\delta_{\Pi_k,k}=\beta_k<1$, whereas
$14k^{k+2}/(3k!)>1$ for $k\ge3$. Thus the same strict global bound holds.
This proves the bound in Theorem~\ref{thm:main}(i); its fixed-$k$ exponent
claim follows from Beniamini--Lavee--Linial's lower bound
\cite[Theorem~7]{BLL}.

\begin{corollary}[absolute-density consequences]\label{cor:growing-coefficient}
Uniformly over $\tau\in\Ssym_k$ and $2\le k=o(N^{1/2})$,
\begin{equation}\label{eq:growing-k-absolute}
 \left|\frac{\#_\tau(\Pi_N)}{\binom Nk}-\frac1{k!}\right|
 \le\left(\frac{14}{3}+O\left(\frac{k^2}{N}\right)\right)\frac{k^3}{N},
\end{equation}
with an absolute implied constant. For $N\ge2k^2+1$, the left side is
less than $7k^3/N$.
\end{corollary}
\begin{proof}
The product $f_{N,k}=\prod_{j<k}(1-j/N)$ is $1+O(k^2/N)$ in the stated
range. Divide Theorem~\ref{thm:main}(i) by
$\binom Nk=(N^k/k!)f_{N,k}$. For the finite bound use
$f_{N,k}\ge1-k(k-1)/(2N)>3/4$ and $56/9<7$.
\end{proof}

\section{Exact rectangular profiles and consequences}\label{sec:rectangles}

The exact binomial-basis formula determines the growing-pattern threshold
and then the local leading vector for fixed $k$.

\subsection{An exact formula in the binomial basis}

The preceding defect argument can be sharpened to an exact formula.  Write
$\alpha\vDash k$ for a positive composition of $k$, let
$\ell(\alpha)$ be its number of parts, and let $p_\alpha(i)$ be the
index of the part containing position $i$.  Thus a composition records
consecutive equality blocks in an ordered list.

Throughout this section the coordinate sets are $[a]$ and $[b]$, so both
coordinates are positive. For $\tau\in\Ssym_k$, compositions
$\alpha,\beta\vDash k$, and $c\in\{-1,+1\}^k$, put
$p_i=(p_\alpha(i),p_\beta(\tau(i)),c_i)$, with block indices starting at one.
Let $w_\tau(\alpha,\beta)$ count the colour words for which
$p_1,\ldots,p_k$ are strictly increasing in $K_X$ and
$p_{\tau^{-1}(1)},\ldots,p_{\tau^{-1}(k)}$ are strictly increasing in $K_Y$.
Only comparisons within equal-coordinate blocks require checking.

For example, if $k=2$, $\alpha=(2)$, $\beta=(1,1)$, and $\tau=12$,
the $Y$-blocks are singletons and the admissible words are $(+,-)$ and
$(-,-)$; hence $w_{12}((2),(1,1))=2$.

\begin{theorem}[exact rectangular profile]\label{thm:exact-rectangle}
For every $a,b\geq1$, every $2\leq k\leq2ab$, and every
$\tau\in\Ssym_k$,
\begin{equation}\label{eq:exact-rectangle}
 \#_\tau(\mathrm{ES}^{\pm}(a,b))
 =\sum_{\alpha,\beta\vDash k}
   w_\tau(\alpha,\beta)
   \binom a{\ell(\alpha)}\binom b{\ell(\beta)}.
\end{equation}
Thus every rectangular profile entry is an explicit bivariate polynomial
in $a$ and $b$.
\end{theorem}

\begin{proof}
List a selected set as $p_1,\ldots,p_k$ in final $X$-order.  Its equal
original $x$-coordinates form consecutive blocks, giving a unique
composition $\alpha$; the equal original $y$-coordinates in final
$Y$-order similarly give $\beta$.  If the induced pattern is $\tau$,
then $p_i$ lies in original-coordinate blocks
$(p_\alpha(i),p_\beta(\tau(i)))$.

Choose the distinct original coordinate values in
$\binom a{\ell(\alpha)}\binom b{\ell(\beta)}$ ways.  Within an equal-$x$
block, the key $-cy$ orders all positive colours with decreasing
$y$-block index, followed by all negative colours with increasing index;
this is exactly the required $K_X$ order. The key $cx$ gives the
required $K_Y$ order symmetrically.  Strictness also excludes selecting
the same coloured point twice.  Hence admissible colour words are in
bijection with the selected sets having these data, proving
\eqref{eq:exact-rectangle}.

\end{proof}

\subsection{The sharp growing-pattern range}

The target $1/k!$ is too small for an absolute collision-probability bound
to give relative balance. We instead retain the exact cancellation of
disjoint isolated pairs before estimating connected interactions.

\begin{lemma}[connected collision bound]\label{lem:connected-collisions}
Put $s=\min(a,b)$ and $n=2ab$. For $2\le k$ and $64k<s$, uniformly over
$\tau\in\Ssym_k$,
\begin{equation}\label{eq:connected-bound}
 \left|\frac{k!\,\#_\tau(\mathrm{ES}^{\pm}(a,b))}{\binom nk}-1\right|
 \le\exp\left(\frac{4096k^3}{s^2(1-64k/s)}
              +\frac{k(k-1)}{2(n-k+1)}\right)-1.
\end{equation}
\end{lemma}
\begin{proof}
Fix $\tau$ and set $\rho=\tau^{-1}$. Form a multigraph on $[k]$ with
$X$-edges $(i,i+1)$ and $Y$-edges $(\rho(j),\rho(j+1))$, retaining
parallel edges. A selected edge set $S$ specifies equality compositions
$\alpha,\beta$ on the two axes. Put
\[
 F(S)=\prod_i\alpha_i!\prod_j\beta_j!,\qquad
 W(S)=2^{-k}w_\tau(\alpha,\beta)F(S).
\]
Let $\mathsf S$ be the equality mask of two independent sorted samples of
$k$ uniform variables from $[a]$ and $[b]$. A composition $\alpha$ on the
first axis has probability $\binom a{\ell(\alpha)}k!/(a^k\prod_i\alpha_i!)$.
Thus \eqref{eq:exact-rectangle} gives the exact identity
\begin{equation}\label{eq:collision-expectation}
 \mathbb EW(\mathsf S)=\frac{(k!)^2\#_\tau(\mathrm{ES}^{\pm}(a,b))}{n^k}.
\end{equation}

The weight $W(S)$ factors over the connected components of $S$. Each
constraint uses the endpoint colours and the relative order or equality
of their opposite-coordinate blocks. That equality forces the connecting
monochromatic path into the same component, so changing blocks outside it
does not alter the constraint. This also covers repeated cells; parallel
edges are one possible connected interaction, not independent factors.
Isolated vertices have weight one. Moreover $W(\varnothing)=W(\{e\})=1$,
by the two-versus-two tie count.

Define $c(T)=\sum_{U\subseteq T}(-1)^{|T|-|U|}W(U)$, so
$W(S)=\sum_{T\subseteq S}c(T)$. Factoring the Boolean sum over components
shows that $c(T)$ is their coefficient product. A component consisting of
one edge therefore makes $c(T)=0$. For a component with $d$ edges,
$0\le W(U)\le F(U)\le F(T)$ for every $U\subseteq T$: splitting a block
can only decrease its factorial product. Consequently
$|c(T)|\le2^dF(T)$, also after multiplying over components.

To control the equality events jointly, realize each discrete sample by
$1+\lfloor aU_i\rfloor$ with independent continuous uniform $U_i$.
An $X$-run of $r$ vertices forces its $r-1$ intervening order-statistic
gaps to have sum at most $1/a$. Distinct runs use disjoint gaps. The joint
density of any $d_X$ gaps is
$(k)_{d_X}(1-\sum z_i)^{k-d_X}\le(k)_{d_X}$ on its simplex. Hence
\[
 \mathbb P(T_X\subseteq\mathsf S)
 \le\frac{(k)_{d_X}a^{-d_X}}{\prod_{X\text{-runs}}(r-1)!},
\]
where $(k)_d=k!/(k-d)!$. This is a joint bound, not an independence
assumption between runs. Applying it independently to the two axes, and
using $r!/(r-1)!=r\le2^{r-1}$, gives for every $T$
\begin{equation}\label{eq:collision-moment}
 |c(T)|\mathbb P(T\subseteq\mathsf S)
 \le(4k/a)^{d_X}(4k/b)^{d_Y}.
\end{equation}

The multigraph has maximum degree four. There are at most $k16^d$ connected
$d$-edge sets: choose a root and encode a tour of length $2d$ traversing
each edge in both directions, with at most four choices per step. Let
$\mathcal C$ contain all connected edge sets of size at least two and put
$q(C)=(4k/a)^{d_X(C)}(4k/b)^{d_Y(C)}$. Expanding into components, applying
\eqref{eq:collision-moment}, and then dropping the disjointness restriction
on component collections yields
\[
 |\mathbb EW-1|\le\prod_{C\in\mathcal C}(1+q(C))-1
 \le\exp\left(k\sum_{d\ge2}(64k/s)^d\right)-1.
\]
Finally, if $f_{n,k}=(n)_k/n^k$, then
$\log f_{n,k}^{-1}\le k(k-1)/(2(n-k+1))$.
Divide \eqref{eq:collision-expectation} by $f_{n,k}$ and use
$|\mathbb EW/f_{n,k}-1|\le e^{\sum_Cq(C)}/f_{n,k}-1$ to obtain
\eqref{eq:connected-bound}.
\end{proof}

\begin{corollary}[relative balance for growing patterns]\label{cor:relative-grid}
If $2\le k=k(N)=o(N^{1/3})$, then, uniformly over $\tau\in\Ssym_k$,
\[
 \frac{k!\,\#_\tau(\Pi_N)}{\binom Nk}
 =1+O\left(\frac{k}{\sqrt N}+\frac{k^3}{N}\right),
\]
with an absolute implied constant.
\end{corollary}
\begin{proof}
For $a=\lfloor\sqrt{\lfloor N/2\rfloor}\rfloor$, the construction gives
\begin{equation}\label{eq:square-sandwich}
 \mathrm{ES}^{\pm}(a,a)\preceq\Pi_N
       \preceq\mathrm{ES}^{\pm}(a+3,a+3),
\end{equation}
where $\preceq$ denotes induced containment. At odd sizes the northeast
point is represented by $(a+2,b+1,+1)$, whose two key ranks are last.
Both square sizes differ from $N$ by $O(\sqrt N)$. In the stated range,
their binomial normalizations differ from $\binom Nk$ by factors
$\exp(O(k/\sqrt N))$, while Lemma~\ref{lem:connected-collisions} gives
relative error $O(k^3/N)$ on either square. Monotonicity of pattern counts
in \eqref{eq:square-sandwich} proves the assertion.
\end{proof}

\begin{proposition}[increasing patterns at the critical scale]
\label{prop:increasing-threshold}
For positive integers $a,b$ and $1\le k\le2ab$, write
$I_{a,b}(k)=\#_{12\cdots k}(\mathrm{ES}^{\pm}(a,b))$. Then
\begin{equation}\label{eq:increasing-exact}
 I_{a,b}(k)=2\sum_{j=0}^{k-1}\binom{k-1}{j}
       \binom{a+k-1-j}{k}\binom{b+j}{k},
\end{equation}
where a binomial coefficient is zero when its upper argument is smaller
than its lower argument. In particular, the longest increasing subsequence
has length $a+b-1$. For the universal family and $2\le k=k(N)\le N$, if $k/N^{1/3}\to c\in[0,\infty)$,
then
\begin{equation}\label{eq:increasing-critical}
 \frac{k!\,\#_{12\cdots k}(\Pi_N)}{\binom Nk}\longrightarrow e^{-c^3/6}.
\end{equation}
If $k/N^{1/3}\to\infty$, the same relative density tends to zero.
\end{proposition}
\begin{proof}
An increasing occurrence is a chain of distinct cells with both coordinates
nondecreasing. A repeated cell is impossible because its two colours have
opposite $X$- and $Y$-orders. Comparing the keys shows that a $+$ colour
at the next point requires a strict increase in $x$, whereas a $-$ colour
requires a strict increase in $y$; these are the only conditions. A colour
word with $j$ plus signs after its first position therefore has
$\binom{a+k-1-j}{k}\binom{b+j}{k}$ possible cell chains. Indeed, a weakly
increasing $k$-tuple in $[a]$ with $j$ specified strict gaps is counted by
$\binom{a+k-1-j}{k}$, by subtracting the number of preceding strict gaps.
Summing over the $2\binom{k-1}{j}$ colour words proves
\eqref{eq:increasing-exact}. Its nonzero terms satisfy
$\max(0,k-b)\le j\le\min(k-1,a-1)$, giving the chain-length assertion.

For a square let $J\sim\operatorname{Bin}(k-1,1/2)$ and
$D_{m,k}=\prod_{r=0}^{k-1}(1-r/(2m^2))$. Reversing one falling factorial in
\eqref{eq:increasing-exact} gives
\begin{equation}\label{eq:increasing-product}
 \frac{k!I_{m,m}(k)}{\binom{2m^2}{k}}
 =\frac1{D_{m,k}}\mathbb E\prod_{r=0}^{k-1}
       \left(1-\frac{(J-r)^2}{m^2}\right).
\end{equation}
This identity holds throughout $1\le k\le2m^2$. Terms outside
$\max(0,k-m)\le J\le\min(k-1,m-1)$ vanish; all factors are positive on this
support. If $k/(2m^2)^{1/3}\to c<\infty$, then $k=o(m)$ and
\[
 \sum_{r=0}^{k-1}(J-r)^2
 =k\left(J-\frac{k-1}{2}\right)^2+\frac{k(k^2-1)}{12}.
\]
After division by $m^2$, the first term has expectation
$k(k-1)/(4m^2)=o(1)$ and the second tends to $c^3/6$.
The uniform logarithmic remainder is $O(k^5/m^4)=o(1)$, while
$\log D_{m,k}=O(k^2/m^2)=o(1)$. The products lie in $[0,1]$, so convergence
in probability also gives convergence of expectations. This proves the
critical limit on squares, including $c=0$.

For $2\le k\le2m-1$, bounding the product in
\eqref{eq:increasing-product} by its quadratic exponential and using
$D_{m,k}^{-1}\le\exp(k(k-1)/(2m^2))$ gives
\[
 \frac{k!I_{m,m}(k)}{\binom{2m^2}{k}}
 \le \exp\left(-\frac{k(k-1)(k-5)}{12m^2}\right).
\]
For $k\ge2m$ the count is zero. Thus the relative increasing density tends
to zero whenever $k^3/m^2\to\infty$.

Finally use \eqref{eq:square-sandwich}. For $k=O(\sqrt N)$, its
square binomial normalizations differ from $\binom Nk$ by
$\exp(O(k/\sqrt N))$. At the critical scale this factor tends to one,
and the square limits squeeze the middle count. In the supercritical
regime the upper-square count is zero for $k\ge2(a+3)$; otherwise the
normalization factor is bounded and the square exponential tends to zero.
Together with Corollary~\ref{cor:relative-grid}, this proves the equivalence
in Theorem~\ref{thm:main}(ii): if $k/N^{1/3}$ does not tend to zero, take
a subsequence converging to a positive finite value or to infinity.
The increasing pattern then prevents relative balance.
\end{proof}

\subsection{The local defect-two formula}

For fixed $k$, the first nonuniform terms have coordinate defect two.
The following local statistics encode their colour sums. Put
$\rho=\tau^{-1}$ and define
\[
\begin{aligned}
 T_x(\tau)&=\#\{i\leq k-2:\tau(i+1)\neq
       \max(\tau(i),\tau(i+1),\tau(i+2))\},\\
 T_y(\tau)&=T_x(\tau^{-1}),\qquad Q_k=\binom{k-2}{2},\\
 H_i&={\bf1}_{i<k,\,\tau(i+1)<\tau(i)}
      -{\bf1}_{i>1,\,\tau(i-1)<\tau(i)},\\
 V_i&={\bf1}_{\tau(i)<k,\,\rho(\tau(i)+1)<i}
      -{\bf1}_{\tau(i)>1,\,\rho(\tau(i)-1)<i},\\
 M(\tau)&=(k-1)^2-\sum_{i=1}^kH_iV_i.
\end{aligned}
\]
Here and below a false subscript condition contributes zero.

\begin{lemma}[explicit defect two]\label{lem:defect-two}
Let $D^{(2)}_\tau(a,b)$ count core subsets of coordinate defect exactly
$2$.  Then
\[
\begin{aligned}
 D^{(2)}_\tau(a,b)=2^{k-2}\bigg[&
  (Q_k+T_x(\tau))\binom a{k-2}\binom bk\\
 &+M(\tau)\binom a{k-1}\binom b{k-1}\\
 &+(Q_k+T_y(\tau))\binom ak\binom b{k-2}\bigg].
\end{aligned}
\]
\end{lemma}

\begin{proof}
Writing $d_x=k-r_x$ and $d_y=k-r_y$, defect two means
$d_x+d_y=2$.  Thus $(d_x,d_y)$ is $(2,0)$, $(1,1)$, or $(0,2)$;
a defect-two equality partition on one axis is either one triple or two
disjoint pairs.  After factoring out $2^{k-2}$ compatible colour choices, the five
 defect-two types and their normalized contributions are recorded in the next display.
\[
\begin{array}{c|c|c}
\text{type}&\text{normalized placements}&\text{coordinate choices}\\ \hline
\text{two disjoint $x$-pairs}&Q_k&\binom a{k-2}\binom bk\\
\text{one $x$-triple}&T_x(\tau)&\binom a{k-2}\binom bk\\
\text{one pair on each axis}&M(\tau)&\binom a{k-1}\binom b{k-1}\\
\text{two disjoint $y$-pairs}&Q_k&\binom ak\binom b{k-2}\\
\text{one $y$-triple}&T_y(\tau)&\binom ak\binom b{k-2}
\end{array}
\]
Indeed, two disjoint $x$-pairs occupy two nonadjacent edges of the
position path, giving $Q_k=\binom{k-2}{2}$.  For an $x$-triple, the
colours form a positive prefix with decreasing ordinates followed by a
negative suffix with increasing ordinates.  If the middle ordinate is
largest, no cut works; otherwise exactly two cuts work.  This gives the
$T_x$ row, and inversion gives the two $y$-rows.

It remains to verify the mixed row.  Let $e=(u,u+1)$ be an adjacent
position edge and let $f=(s,s+1)$ be an adjacent value edge, pulled back to
positions $\rho(s),\rho(s+1)$.  Write their active-endpoint functions as
\[
 h_e(i)={\bf1}_{\{i=u,\,\tau(u)>\tau(u+1)\}}
        -{\bf1}_{\{i=u+1,\,\tau(u)<\tau(u+1)\}},
\]
\[
 v_f(i)={\bf1}_{\{i=\rho(s),\,\rho(s)>\rho(s+1)\}}
        -{\bf1}_{\{i=\rho(s+1),\,\rho(s+1)>\rho(s)\}}.
\]
First suppose that $e$ and the pulled-back
$f$ have the same two endpoints.  The two selected labels then lie over one
cell, so distinctness forces their colours to be opposite.  Comparing the two
key orders shows that the $+$ point comes first in $X$ and second in $Y$.
Consequently the induced pair is necessarily a descent: the normalized weight
is one for the descent and zero for the ascent.  In the descent the two active
endpoints below are different, so $\sum_i h_e(i)v_f(i)=0$; in the ascent they
coincide with product one.  Thus the desired weight
$1-\sum_i h_e(i)v_f(i)$ already holds in this repeated-cell case.

Assume now that the two equality edges do not have the same pair of endpoints.
Each equality imposes exactly one colour constraint.  If
$\tau(u)>\tau(u+1)$, the two points on $e$ appear in the required $X$-order
exactly when the left endpoint has colour $+$; if
$\tau(u)<\tau(u+1)$, exactly when the right endpoint has colour $-$.  Thus the
active endpoint of $e$ and its forced colour are encoded by the unique nonzero
value $h_e(i)\in\{1,-1\}$.  For $f$, the active endpoint is the later of
$\rho(s),\rho(s+1)$, and its forced colour is $-v_f(i)$.  These assertions
follow immediately by comparing the tied keys $-cy$ and $cx$, respectively.

If the two active endpoints are different, the constraints concern two
colours, so there are $2^{k-2}$ compatible words and the normalized weight is
one.  If they coincide at $i$, two cases remain:
\[
\begin{array}{c|c|c}
 \text{relation}&\text{two constraints}&
 \text{normalized compatible-word count}\\ \hline
 h_e(i)=-v_f(i)&\text{agree}&2\\
 h_e(i)= v_f(i)&\text{conflict}&0.
\end{array}
\]
Both cases are $1-h_e(i)v_f(i)$.  Because each edge has only one active
endpoint, summing over all common endpoints gives the normalized weight
\[
 1-\sum_i h_e(i)v_f(i)
\]
for every noncoincident pair $(e,f)$.  Together with the repeated-cell
calculation above, this proves the formula for every edge pair.  Summing over
all $(k-1)^2$ pairs gives
\[
 (k-1)^2-
 \sum_i\left(\sum_{e\ni i}h_e(i)\right)
       \left(\sum_{f\ni i}v_f(i)\right)
 =(k-1)^2-\sum_iH_iV_i=M(\tau).
\]
The coordinate choices in the table now prove the lemma.
\end{proof}

Define the centered local vector and its norm by
\begin{align}
 g_k(\tau)&=T_x(\tau)+T_y(\tau)+\frac{11-4k}{3}
 -\frac{k}{k-1}\sum_{i=1}^kH_iV_i,\label{eq:g-local}\\
 \lambda_{k,\tau}&=\frac{2^{k-2}}{k!(k-2)!}g_k(\tau),\qquad
 \Lambda_k=\frac{2^{k-2}}{k!(k-2)!}
       \max_\tau|g_k(\tau)|.\label{eq:lambda-local}
\end{align}
For each fixed $k$,
\begin{align}
 \#_\tau(\mathrm{ES}^{\pm}(m,m))
 -\frac1{k!}\binom{2m^2}{k}
 &=\lambda_{k,\tau}m^{2k-2}+O_k(m^{2k-3}),\label{eq:leading-pattern}\\
 \delta_{\mathrm{ES}^{\pm}(m,m),k}
 &=\Lambda_km^{2k-2}+O_k(m^{2k-3}).\label{eq:leading-delta}
\end{align}

These statistics compress the two one-axis triple classes and the mixed
pair class in the preceding table.  For example, $\tau=4213$ has
$T_x=2$, $T_y=1$, and $\sum_iH_iV_i=0$, so
$g_4(4213)=4/3$ and $\lambda_{4,4213}=1/9$.
In the square case, the coefficient in Lemma~\ref{lem:defect-two}, after
factoring out $2^{k-2}/(k!(k-2)!)$, is
\[
 R_k(\tau)=T_x(\tau)+T_y(\tau)+2Q_k+
             \frac{k}{k-1}M(\tau).
\]
For uniform $\tau\in\Ssym_k$, write $\gamma_k=2^{k-2}/(k!(k-2)!)$. The defect-zero and defect-one
classes are pattern-uniform and together contribute
\[
 P_{01}(m)=2^k\left[\binom mk^2+(k-1)\binom mk\binom m{k-1}\right].
\]
There are $k-1$ possible doubled blocks on each axis, each with colour
weight $2^{k-1}$, which gives the second term. Expanding the falling
factorials, the coefficient of $m^{2k-2}$ in $P_{01}$ is
$\gamma_k(-6k^2+14k-10)/3$. The profile average is
$\binom{2m^2}{k}/k!$, whose corresponding coefficient is $-\gamma_k$.
Since the defect-two contribution is $\gamma_kR_k(\tau)$, it follows that
\[
 \mathbb ER_k=\frac{6k^2-14k+7}{3},\qquad
 R_k(\tau)-\mathbb ER_k
 =T_x+T_y+\frac{11-4k}{3}-\frac{k}{k-1}\sum_iH_iV_i.
\]
This identifies the centered defect-two coefficient as $\lambda_{k,\tau}$.
All higher defects have degree at most $2k-3$, proving
\eqref{eq:leading-pattern}; taking the maximum over the fixed finite set
$\Ssym_k$ gives \eqref{eq:leading-delta}. Also
$\mathbb ET_x=\mathbb ET_y=2(k-2)/3$, since each triple has a nonmaximal
middle entry with probability $2/3$. Substitution in $R_k$ gives
$\mathbb E\sum_iH_iV_i=(k-1)/k$ and hence
$\mathbb EM=(k-1)^2-(k-1)/k$, including $k=2$.
After forming
$\tau^{-1}$, each $g_k(\tau)$ costs $O(k)$ to evaluate, so
$\Lambda_k$ is computable in $O(k\,k!)$ time.

\subsection{The fixed-aspect formula and a four-point baseline}
The exact formula and the defect-two lemma also give the following
unoptimized fixed-aspect consequence.  Put
$\mu_k=2(k-2)/3$ and $S(\tau)=\sum_iH_iV_i$.  For $\theta>0$, define
\begin{equation}\label{eq:h-local}
 h_{k,\tau}(\theta)
 =\theta^{-1}(T_x(\tau)-\mu_k)
  +1-\frac{k}{k-1}S(\tau)
  +\theta(T_y(\tau)-\mu_k).
\end{equation}
If $a,b\to\infty$ with $a/b\to\theta$, the defect-two terms in
Lemma~\ref{lem:defect-two}, after centering by the uniform profile, give
\begin{align}
 \#_\tau(\mathrm{ES}^{\pm}(a,b))
 -\frac1{k!}\binom{2ab}{k}
 &=\frac{2^{k-2}(ab)^{k-1}}{k!(k-2)!}
   h_{k,\tau}(a/b)+O_k((a+b)^{2k-3}),
 \label{eq:rectangular-leading}\\
 \frac{\#_\tau(\mathrm{ES}^{\pm}(a,b))-\binom{2ab}{k}/k!}{(2ab)^{k-1}}
 &=\frac{h_{k,\tau}(a/b)}{2k!(k-2)!}+o_k(1).
 \label{eq:rectangular-leading-normalized}
\end{align}
Indeed, the $(2,0)$, $(1,1)$, and $(0,2)$ defect classes have respectively
aspect factors $\theta^{-1}$, $1$, and $\theta$; their uniform averages are
$\mu_k$, $1-\frac{k}{k-1}\mathbb E S$, and $\mu_k$, as computed above.
All higher-defect terms have degree at most $2k-3$.  In particular,
$h_{k,\tau}(1)=g_k(\tau)$, and exchanging the axes exchanges $\tau$ with
$\tau^{-1}$ and $\theta$ with $\theta^{-1}$.  The aspect weights separate the two one-axis triple contributions from
the mixed contribution.  Section~\ref{sec:sharp-four} uses this separation
to optimize the residual after boundary completion and insertions.

For $k=4$, one direct local evaluation is enough for the comparison needed
here:
\[
 h_{4,1243}(\theta)=-\frac{\theta^{-1}+1+\theta}{3},
 \qquad
 \frac{|h_{4,1243}(\theta)|}{96}
 =\frac{\theta^{-1}+1+\theta}{288}\geq\frac1{96}>\frac1{192}.
\]
Thus the oriented inserted family of Theorem~\ref{thm:sharp-four} improves
every standard fixed-aspect complete core.  The square core has the exact
formulas below.

\subsection{Exact four-point profiles and the published square coefficient}
For $k=4$, direct collection of the composition terms in
\eqref{eq:exact-rectangle} gives six exact polynomial classes for $m\ge2$
(some coincide at small $m$):
\[
\begin{array}{c|c|l}
\displaystyle
 \frac{144}{m^2(m^2-1)}
 \left(\#_\tau-\frac1{24}\binom{2m^2}{4}\right)
 &\#&\tau\\ \hline
 -5(4m^2-9)&1&1234\\
 -(12m^2-13)&3&1243,1324,1432\\
 4m^2-3&9&\begin{gathered}1342,1423,2143,2413,2431,\\
                         3142,3412,4132,4231\end{gathered}\\
 -(4m^2-5)&3&2134,3214,4321\\
 -11&6&2314,2341,3124,3421,4123,4312\\
 16m^2-3&2&3241,4213
\end{array}
\]
The table is an exact substitution in the finite rectangle formula, with no
asymptotic truncation.  It yields
\begin{equation}\label{eq:k4-exact}
\delta_{\mathrm{ES}^{\pm}(m,m),4}
=\begin{cases}
 \dfrac{m^2(m^2-1)(16m^2-3)}{144},&m=2,3,\\[4pt]
 \dfrac{5m^2(m^2-1)(4m^2-9)}{144},&m\ge4.
\end{cases}
\end{equation}
In particular the square-grid coefficient is $\Lambda_4=5/36$ against
$m^6$, or $5/288$ against $N^3$ since $N=2m^2$.  For every $k\ge2$, \eqref{eq:g-local} gives
$g_k(12\cdots k)=-(k+1)/3$, so $\Lambda_k>0$.

\begin{remark}[comparison with the published square coefficient]
\label{rem:bll-correction}
In \eqref{eq:keys}, colours $c=+1$ and $c=-1$ correspond respectively
to the rotated copies $P(\varepsilon)$ and $P(-\varepsilon)$ of
\cite[Definitions~5.3 and 5.6]{BLL}; Section~\ref{sec:construction}
verifies equality of both coordinate orders.
Their square side $n$ is our $m$: both constructions have $N=2m^2$
and use the centred profile statistic.  Our exact formula therefore
gives $5m^6/36=5N^3/288$ as the leading term, rather than the printed
$(2/9)m^6$ in
\cite[Theorem~8]{BLL}.  This bibliographic correction is not used in any
bound: the coefficient follows here from
Theorem~\ref{thm:exact-rectangle}, Lemma~\ref{lem:defect-two}, and
\eqref{eq:k4-exact}.  For instance, at $m=4$ the exact count is
$\#_{1234}(\mathrm{ES}^{\pm}(4,4))=1040$.
\end{remark}

\section{Bounded insertions and fixed-\texorpdfstring{$k$}{k} realization}\label{sec:sharp-four}
The local core vector records what a bounded-prefix row and finitely many
insertions can change.  For standard sheets, a moment realization attains
the resulting LP bound with a nonconstructive family
$\Xi_N^{(k,\theta)}$ for fixed $k$ and aspect.  A separate periodic orientation gives the explicit
$k=4$ family $\Sigma_N$, with $12$/$17$ insertions and constant $1/192$.
Both optima are relative to their stated completion classes and neither
replaces the universal family $\Pi_N$.

\subsection{The boundary leading vector}

For a fixed $0<\delta<1$ and integers $0<r<b$, define the shifted
mechanical set
\[
 R_\delta(b,r)=
 \left\{\left\lfloor\frac{(j-\delta)b}{r}\right\rfloor+1:
          1\le j\le r\right\}\subseteq[b].
\]
Its elements are distinct, and for every $0\le z\le b$,
\begin{equation}\label{eq:shifted-prefix}
 \left|\,|R_\delta(b,r)\cap[z]|-\frac rbz\right|\le2.
\end{equation}
Indeed, with $y_j=\lfloor(j-\delta)b/r\rfloor+1$, the condition
$y_j\le z$ is $j<\delta+rz/b$; the number of eligible integers differs
from $rz/b$ by at most two.  Thus every shifted selector satisfies the
hypothesis of Proposition~\ref{prop:boundary-leading}.  The midpoint set
$R_{1/2}(b,r)$ is generally distinct from the ceiling selector $R_{b,r}$
used in Section~\ref{sec:construction}.

For $1\le t\le k$ write
\[
 f_{k,t}^{(b)}(y)=
 \binom{y-1}{t-1}\binom{b-y}{k-t},
\]
and, for $R\subseteq[b]$ with $|R|=r$, put
\begin{equation}\label{eq:q-boundary}
 q_{k,t}(b,R)=\frac1{b^{k-1}}
 \left(\sum_{y\in R}f_{k,t}^{(b)}(y)
       -\frac rb\binom bk\right).
\end{equation}
This is the normalized error from replacing $R$ by density $r/b$ against
this rank weight.  For example, $R_{1/2}(5,2)=\{2,4\}$ gives
$q_{4,3}=1/125$.  At scale $N^{k-1}$ it is the only pattern-dependent
boundary term; collision terms become uniform after density replacement.

\begin{lemma}[bounded-prefix Abel estimate]\label{lem:bounded-prefix-abel}
Let $g$ be supported on an integer interval $I=[u,v]$, and suppose every
initial sum on that interval satisfies
\[
 \left|\sum_{y=u}^z g(y)\right|\le D
 \qquad (u\le z\le v).
\]
Then every weight $W:I\to\mathbb R$ satisfies
\begin{equation}\label{eq:bounded-prefix-abel}
 \left|\sum_{y=u}^v g(y)W(y)\right|
 \le D\left(|W(v)|+
       \sum_{y=u}^{v-1}|W(y+1)-W(y)|\right).
\end{equation}
If $\sum_I g=0$, the endpoint term $D|W(v)|$ may be omitted.  In
particular, if $D=O(1)$, $W$ is a polynomial of degree at most $d$, and
$\lVert W\rVert_\infty=O_d(b^d)$ on $I\subseteq[b]$, then the left side is
$O_{d,D}(b^d)$, uniformly in the subinterval endpoints.  These hypotheses
hold for the rank-completion weights below, whose binomial factors have
arguments between $0$ and $b$.
\end{lemma}

\begin{proof}
Set $G(z)=\sum_{y=u}^z g(y)$.  Summation by parts gives
$\sum gW=G(v)W(v)+\sum_{y=u}^{v-1}G(y)(W(y)-W(y+1))$,
which proves \eqref{eq:bounded-prefix-abel}; $G(v)=0$ removes the endpoint
term.  A fixed-degree polynomial has $O_d(1)$ monotonicity intervals, so
its total variation is $O_d(b^d)$ when its size is $O_d(b^d)$.  The same
size bound holds for the rank-completion weights because a degree-$j$
binomial factor is at most $b^j/j!$.
\end{proof}

\begin{proposition}[boundary leading vector]\label{prop:boundary-leading}
Fix $k\ge2$.  Suppose $a,b\to\infty$, $a/b\to\theta>0$, and
$R\subseteq[b]$, $|R|=r$, satisfies
\[
 \sup_{0\le z\le b}
 \left|\,|R\cap[z]|-\frac rbz\right|=O(1).
\]
For
\[
 F=([a]\times[b])\cup(\{a+1\}\times R),
 \qquad N=2|F|,
\]
one has, uniformly over $\tau\in\Ssym_k$,
\begin{equation}\label{eq:boundary-leading}
 \frac{\#_\tau(\Pi(F))-\binom Nk/k!}{N^{k-1}}
 =\frac{h_{k,\tau}(\theta)}{2k!(k-2)!}
  +\frac{2}{(k-1)!}q_{k,\tau(k)}(b,R)+o(1).
\end{equation}
\end{proposition}

\begin{proof}
Let $C$ be a fixed bound for the prefix discrepancy.  Use the composition--colour-word
partition from Theorem~\ref{thm:exact-rectangle}.  For
$\alpha=(\alpha_1,\ldots,\alpha_p)$ put
\[
 I_\alpha=\{k-\alpha_p+1,\ldots,k\},
\]
the last $X$-block.  When a term uses the partial row, this block has
abscissa $a+1$.  For $\alpha,\beta\vDash k$, set
\[
 Q_{\alpha,\beta,\tau}
 =\{p_\beta(\tau(i)):i\in I_\alpha\}\subseteq[\ell(\beta)]
\]
and, for $Q\subseteq[q]$, set
\[
 S_Q(b,R)=\sum_{1\le y_1<\cdots<y_q\le b}
               \prod_{j\in Q}{\bf1}_R(y_j).
\]
The partial-row contribution of $(\alpha,\beta)$ is exactly
\begin{equation}\label{eq:partial-composition}
 \binom a{p-1}w_\tau(\alpha,\beta)
 S_{Q_{\alpha,\beta,\tau}}(b,R),
 \qquad p=\ell(\alpha),\quad q=\ell(\beta).
\end{equation}
Indeed, the other $p-1$ abscissae lie in $[a]$, and precisely the ordinate
levels meeting the last $X$-block must lie in $R$.  This condition depends
only on the coordinate blocks, so the full colour sum $w_\tau$ is retained.

The summand in \eqref{eq:partial-composition} is
$O_k(a^{p-1}b^q)$.  Therefore $p+q\le2k-2$ contributes
$O_k(b^{2k-3})=o(N^{k-1})$.  For the partial-row part with $p=q=k$,
$w_\tau(1^k,1^k)=2^k$ and $Q=\{\tau(k)\}$, giving
\[
 2^k\binom a{k-1}\sum_{y\in R}f_{k,\tau(k)}^{(b)}(y).
\]
The profile average of the sum is $(r/b)\binom bk$ by Vandermonde;
after division by $N^{k-1}$, the centred remainder is the
second term of \eqref{eq:boundary-leading}.

It remains to handle $p+q=2k-1$, for which $|Q|\le2$.  Uniformly for fixed
$Q\subseteq[q]$ with $|Q|\le2$,
\begin{equation}\label{eq:density-replacement}
 S_Q(b,R)=\left(\frac rb\right)^{|Q|}\binom bq
          +O_{k,C}(b^{q-1}).
\end{equation}
For $|Q|=0$ this is exact.  For $|Q|=1$, apply
Lemma~\ref{lem:bounded-prefix-abel} to $g={\bf1}_R-r/b$, whose prefix sums
are bounded by $C$, and to
$W_{q,j}(y)=\binom{y-1}{j-1}\binom{b-y}{q-j}$; its degree is $q-1$, so
the error is $O_{k,C}(b^{q-1})$.  If $Q=\{j,\ell\}$ with $j<\ell$,
expand each indicator as $r/b+g$.  The one-$g$ terms are covered by the
same argument.  For the two-$g$ term, fix $v=y_\ell$; the inner sum over
$u=y_j<v$ has weight
\[
 U_v(u)=\binom{u-1}{j-1}
        \binom{v-u-1}{\ell-j-1}
        \binom{b-v}{q-\ell},
\]
with inadmissible binomial coefficients zero.  Every interval sum of $g$
has absolute value at most $2C$.  The support of $U_v$ is
$\{j,\ldots,v-(\ell-j)\}$, and on it $U_v$ has degree at most $q-2$,
supremum $O_k(b^{q-2})$, and variation $O_k(b^{q-2})$.  Abel summation,
retaining its endpoint term, gives uniformly in $v$
\[
 \left|\sum_{u<v}g(u)U_v(u)\right|
 \le 2C\bigl(\lVert U_v\rVert_\infty+\operatorname{Var}(U_v)\bigr)
 =O_{k,C}(b^{q-2}).
\]
Since $|g(v)|\le1$, summing over $v$ gives $O_{k,C}(b^{q-1})$ and proves
\eqref{eq:density-replacement}.

The only leading partial-row types are $(p,q)=(k-1,k)$ and $(k,k-1)$.
In the first, $\beta=1^k$, $\alpha$ has one part of size two, and
$|Q|=\alpha_p\in\{1,2\}$, independently of $\tau$; in the second,
$\alpha=1^k$ and $|Q|=1$.  After
\eqref{eq:density-replacement}, the first main term is
$(r/b)^{\alpha_p}$ times the corresponding complete-row class with the
largest abscissa fixed at $a+1$.  For each fixed original $x$-multiset and
set of distinct ordinates, the first part of Lemma~\ref{lem:core-uniform},
with its full colour sum, makes this vector pattern-uniform.  The second
main term is $r/b$ times the analogous class with distinct abscissae and one
repeated ordinate, which is uniform by the symmetric part of that lemma.  No
claim is needed for an individual colour word; \eqref{eq:partial-composition}
retains the complete sum, and the scalar depends only on the composition.

Each replacement error is
$O_{k,C}(a^{p-1}b^{q-1})=O_{k,C}(b^{2k-3})$, and there are finitely many
compositions for fixed $k$.  Thus all partial-row main terms other than the
clean boundary term are pattern-uniform and all errors are
$o(N^{k-1})$.  The complete core contributes the first term of
\eqref{eq:boundary-leading} by \eqref{eq:rectangular-leading}.
\end{proof}

\begin{corollary}[universal-family coefficient]\label{cor:fixed-coefficient}
For every fixed $k\ge2$, with $\Lambda_k$ as in
Theorem~\ref{thm:exact-rectangle},
\[
 \limsup_{N\to\infty}\frac{\delta_{\Pi_N,k}}{N^{k-1}}
 \le\frac{\Lambda_k}{2^{k-1}}+\frac{3k-1}{k!(k-1)!}.
\]
\end{corollary}
\begin{proof}
Lemma~\ref{lem:rank-balance} gives $|q_{k,t}(b,R_{b,r})|\le1/(k-1)!$.
Proposition~\ref{prop:boundary-leading}, with $a/b\to1$, bounds the even-order
limsup by $\Lambda_k/2^{k-1}+2/((k-1)!)^2$.
For odd orders, \eqref{eq:padding-sharp} adds
$(k-1)/(k!(k-1)!)$: the term $\delta_{\Pi_{N-1},k-1}$ is lower order for
fixed $k\ge3$, and $\delta_{\Pi_{N-1},1}=0$. Summing proves the claim.
\end{proof}

\begin{lemma}[midpoint quadrature]\label{lem:midpoint-quadrature}
Suppose $b,r\to\infty$ and $r/b\to\lambda\in(0,1)$ is irrational.  For
every fixed $k$ and $t$,
\[
 q_{k,t}\bigl(b,R_{1/2}(b,r)\bigr)\longrightarrow0.
\]
\end{lemma}

\begin{proof}
Let
\[
 A_z=|R_{1/2}(b,r)\cap[z]|-\frac rbz
 \qquad(0\le z\le b).
\]
Vandermonde and summation by parts give
\begin{equation}\label{eq:midpoint-abel}
 q_{k,t}\bigl(b,R_{1/2}(b,r)\bigr)
 =-\frac1{b^{k-1}}\sum_{z=1}^{b-1}
 A_z\bigl(f_{k,t}^{(b)}(z+1)-f_{k,t}^{(b)}(z)\bigr).
\end{equation}
Write $\langle x\rangle=x-\lceil x\rceil+1$.  Then
$A_z=\frac12-\langle\frac12+rz/b\rangle$, and the pairs
\[
 \left(\frac zb,\left(\frac12+\frac{rz}{b}\right)\!\!\pmod1\right),
 \qquad 1\le z<b,
\]
are equidistributed in $[0,1]^2$.  Indeed, for fixed integers $m,n$ not
both zero the Fourier sum is geometric with ratio
$\exp(2\pi i(m+nr)/b)$.  For $n\ne0$ this tends to
$\exp(2\pi i n\lambda)\ne1$, while for $n=0$ the complete roots-of-unity
sum has normalized magnitude $O(1/b)$.  The integer endpoints of
$\langle\cdot\rangle$ account for $o(b)$ indices, since the reduced
denominators tend to infinity.

Put
$g_{k,t}(x)=x^{t-1}(1-x)^{k-t}/((t-1)!(k-t)!)$.  Uniformly for
$0\le z<b$,
\[
 \frac{f_{k,t}^{(b)}(z+1)-f_{k,t}^{(b)}(z)}{b^{k-2}}
 =g'_{k,t}(z/b)+O_k(1/b).
\]
Thus \eqref{eq:midpoint-abel} is an equidistributed Riemann sum for
$(\frac12-u)g'_{k,t}(x)$, whose integral is zero because
$\int_0^1(\frac12-u)\,du=0$.  The uniform remainder and the exceptional
indices are $o(1)$, proving the limit.
\end{proof}

\begin{lemma}[dimension search]\label{lem:dimension-search}
Fix $\lambda\in(0,1)$. For sufficiently large integers $E$, put
$B=\lfloor\sqrt E\rfloor$ and $L=\lceil B^{2/3}\rceil$.
Among $b=B,\ldots,B+L$, choose the smallest minimizer of
$d(E/b,\lambda)$, where $d(u,v)=\min_{m\in\mathbb Z}|u-v-m|$,
and set $a=\lfloor E/b\rfloor$, $r=E-ab$. Then
\begin{equation}\label{eq:dimension-search}
 a/b\longrightarrow1,\qquad r/b\longrightarrow\lambda.
\end{equation}
\end{lemma}
\begin{proof}
Every candidate has $b/B\to1$ and $a/b\to1$. For
$0\le j\le J=\lfloor2\sqrt B\rfloor<L$, put $F_j=E/(B+j)+j$.
Uniformly in this range,
\[
 F_j=E/B+j^2/B+O(B^{-1/2}),\qquad
 F_{j+1}-F_j=O(B^{-1/2}).
\]
Since $F_J-F_0\to4$, the first crossing of a level in
$\mathbb Z+\lambda$ between the endpoints has overshoot $O(B^{-1/2})$.
Subtracting the integer $j$ does not change its distance modulo one,
so the minimum tends to zero. As $\lambda\in(0,1)$, this gives the
ordinary limit $r/b\to\lambda$ and eventually $0<r<b$.
\end{proof}

\subsection{Uniform fixed-\texorpdfstring{$k$}{k} insertion expansion}
A \emph{point insertion} adjoins an element to both defining total orders
without changing the old relative orders.  With $n$ old elements, encode its
slot by normalized ranks $(u,v)=(x/n,y/n)$, where $x,y$ count old elements
preceding it.  All slots are allowed, including endpoints and slots inside
coordinate blocks or the partial row; total orders resolve displayed ties.

For fixed $k$, put
\[
 A_k=2k!(k-2)!,\qquad s_k=((k-1)!)^2,
 \qquad B^{(k)}_j(u)=\binom{k-1}{j}u^j(1-u)^{k-1-j}
 \quad(0\leq j\leq k-1).
\]
The following lemma is the uniform insertion statement used by both the
explicit four-point construction and the fixed-$k$ theorem.

\begin{lemma}[uniform fixed-$k$ insertion expansion]\label{lem:finite-padding}
Fix $k\ge2$.  Suppose $a,b\to\infty$, $a/b\to\theta>0$, and
\[
 F=([a]\times[b])\cup(\{a+1\}\times R),
 \qquad n=2|F|,
\]
where $R\subseteq[b]$ has bounded prefix discrepancy as in
Proposition~\ref{prop:boundary-leading}.  Fix $M$, and make $m\le M$ point
insertions with normalized ranks $(u_j,v_j)\in[0,1]^2$, or with arbitrary
slots having those normalized ranks.  If $N=n+m$ is the resulting size,
then, uniformly over the list, the slots, and $\tau\in\Ssym_k$,
\begin{equation}\label{eq:insertion-leading}
\begin{aligned}
 \frac{\#_\tau(\sigma)-\binom Nk/k!}{N^{k-1}}
 &=\frac{h_{k,\tau}(\theta)}{A_k}
  +\frac{2}{(k-1)!}q_{k,\tau(k)}(b,R)\\[-2pt]
 &\quad+\sum_{j=1}^m\left[
    \frac1{s_k}\sum_{i=1}^k
      B^{(k)}_{i-1}(u_j)B^{(k)}_{\tau(i)-1}(v_j)
    -\frac1{k!(k-1)!}\right]+o(1).
\end{aligned}
\end{equation}
For asymptotic purposes, a slot with normalized ranks $(u_j,v_j)$ may,
when $m\ge1$, be represented by the point
\begin{equation}\label{eq:inserted-points}
 p_j=\left(\lfloor u_ja\rfloor+\frac{j}{m+1},
            \lfloor v_jb\rfloor+\frac{j}{m+1},0\right),
 \qquad 1\le j\le m.
\end{equation}
The fractional offsets separate the displayed inserted points; arbitrary
slots, including endpoints and slots inside coordinate blocks, are handled
by the movement estimate in the proof.  The old points use the same
coordinate keys as in \eqref{eq:keys}.
\end{lemma}

\begin{proof}
For the displayed representatives, configurations containing at least two
inserted points number $O_{k,M}(n^{k-2})$; the case $m=0$ is empty.  Fix
one inserted point $p_j$ and suppose it has position $i$ in the selected
$X$-order and value position $\tau(i)$ in the selected $Y$-order.  After
excluding a partial-row point and all old coordinate collisions, its leading
coordinate-set contribution is
\[
 2^{k-1}\binom{\lfloor u_ja\rfloor}{i-1}
              \binom{a-\lfloor u_ja\rfloor}{k-i}
              \binom{\lfloor v_jb\rfloor}{\tau(i)-1}
              \binom{b-\lfloor v_jb\rfloor}{k-\tau(i)}.
\]
The $O(1)$ fractional offsets in \eqref{eq:inserted-points} are covered by
the uniform binomial estimate below.  The factor $2^{k-1}$ counts the
independent old colours; once the two coordinate sets and the insertion
ranks are fixed, the pattern fixes the cell bijection.

Uniformly for $0\le u\le1$ and fixed $d\ge1$,
\[
 \binom{\lfloor ua\rfloor+O(1)}d
 =\frac{(ua)^d}{d!}+O_d(a^{d-1}),
\]
with the analogous estimate for $b,v$ and exact value one when $d=0$.
The absolute estimates include $u,v=0,1$.  Summing over $i$ and using
$N^{k-1}\sim(2ab)^{k-1}$ gives
\[
 \frac1{s_k}\sum_{i=1}^k
 B^{(k)}_{i-1}(u_j)B^{(k)}_{\tau(i)-1}(v_j)
 +O_k(a^{-1}+b^{-1}).
\]
The omitted repeated-old-coordinate and partial-row configurations total
\[
 O_k(a^{k-2}b^{k-1}+a^{k-1}b^{k-2}+a^{k-2}rb^{k-2})
 =O_{k,\theta}(n^{k-3/2}),
\]
uniformly in the cuts, hence are $o(n^{k-1})$ because $r<b$ and
$a/b\to\theta$.  The same holds for the multi-insertion configurations.
The target changes by
\[
 \frac1{k!}\left(\binom{n+m}{k}-\binom nk\right)
 =\frac{m\,n^{k-1}}{k!(k-1)!}+O_{k,M}(n^{k-2}),
\]
which supplies the bracket in \eqref{eq:insertion-leading}.  Proposition~\ref{prop:boundary-leading}
supplies the core and row terms, while its Abel estimate gives
$q_{k,t}(b,R)=O_k(1)$.  This proves the formula for the representatives,
uniformly in the macro-coordinates.

For an arbitrary slot, a cut at $x=ua$ has $2b$ old points per complete
column and at most $2|R|=O(b)$ partial-row points; its old rank therefore
differs from the prescribed rank by $O(a+b)$.  The analogous $y$-cut bound
uses column masses $2a$ and at most two extra points.  Moving an insertion
through $D=O(a+b)=O(\sqrt n)$ old positions changes at most
$\binom{n+m-2}{k-2}=O_{k,M}(n^{k-2})$ configurations per position; all
configurations with two inserted points total $O_{k,M}(n^{k-2})$.  The
movement error is consequently
$O_{k,M}(n^{k-3/2})=o(n^{k-1})$.  Finally, perturbing finitely many
coincident displayed coordinates by $o(1)$ in normalized units removes all
old and new coordinate ties.  Continuity of the Bernstein factors preserves
the leading term and the endpoint uniformity.
\end{proof}

\subsection{Bounded moment realization}
We realize any bounded centered correction with a bounded number of insertions.
Let
\[
 \mathbf1=(1,\ldots,1)^{\mathsf T},\qquad
 \mathcal P_k=I_k-\frac1k\mathbf1\mathbf1^{\mathsf T},\qquad
 \mathcal V_k=\{T\in\mathbb R^{k\times k}:T\mathbf1=0,\ \mathbf1^{\mathsf T}T=0\}.
\]
For a finite ordered list $Z=((u_\ell,v_\ell))_{\ell=1}^m$ in $[0,1]^2$, put
\[
 \Phi_k(Z)=\mathcal P_k\left(\sum_{\ell=1}^m
 B^{(k)}(u_\ell)B^{(k)}(v_\ell)^{\mathsf T}\right)\mathcal P_k,
 \qquad
 B^{(k)}(u)=(B^{(k)}_0(u),\ldots,B^{(k)}_{k-1}(u))^{\mathsf T}.
\]
Then $\Phi_k(Z)\in\mathcal V_k$.  If
$Q=\sum_{\ell=1}^mB^{(k)}(u_\ell)B^{(k)}(v_\ell)^{\mathsf T}$, then
$\sum_{i,j}Q_{ij}=m$ and, for every $\tau\in\Ssym_k$,
\begin{equation}\label{eq:centered-assignment}
 \sum_{i=1}^k(\mathcal P_kQ\mathcal P_k)_{i,\tau(i)}
 =\sum_{i=1}^kQ_{i,\tau(i)}-\frac mk.
\end{equation}

\begin{lemma}[compact centered moment realization]\label{lem:moment-realization}
For every compact set $K\subset\mathcal V_k$, there are an integer $m$ and a
continuous choice of an $m$-point list $Z(T)$, $T\in K$, with all coordinates
in a fixed compact subset of $(0,1)^2$, such that
\[
 \Phi_k(Z(T))=T\qquad(T\in K).
\]
The integer $m$ may be chosen even.  There is also an odd integer $h$ and an
$h$-point neutral list $Z_{\rm odd}$ with $\Phi_k(Z_{\rm odd})=0$.
\end{lemma}

\begin{proof}
For general averaging sets see Seymour--Zaslavsky~\cite{SeymourZaslavsky}.
Take a positive $k$-node Gauss rule $(x_i,w_i)$ on $[0,1]$, with distinct
interior nodes, exact through degree $2k-1$. For every sufficiently large
$L$, choose positive integers $n_i$ with $\sum_i n_i=L$ and
$n_i/L=w_i+O(1/L)$. Fix one node $x_{i_0}$. The Jacobian
$(d w_i x_i^{d-1})_{1\le d\le k-1,\ i\ne i_0}$ is an invertible
Vandermonde scaling, so the implicit-function theorem adjusts the other
nodes by $O(1/L)$ to restore moments $1/(d+1)$ for $1\le d\le k-1$.
Degree zero is automatic. The resulting equal-weight list $U_L$ satisfies
$\sum_{u\in U_L}B^{(k)}(u)=(L/k)\mathbf1$, with support in a fixed interior
compact interval. This works for every large $L$, including odd $L$.

Fix one such $L$.  The product $Z_0=U_L\times U_L$ is neutral, with
uncentered moment matrix $L^2\mathbf1\mathbf1^{\mathsf T}/k^2$.  Moving the
$u$-coordinate of one indexed copy at each distinct support pair gives
Jacobian columns
$B^{(k)\prime}(x_i)(\mathcal P_kB^{(k)}(y_j))^{\mathsf T}$.  The vectors
$B^{(k)\prime}(x_i)$ span the zero-sum hyperplane: a vector orthogonal to
all of them gives a Bernstein polynomial whose derivative vanishes at $k$
distinct points, hence is constant.  The vectors
$\mathcal P_kB^{(k)}(y_j)$ span the same hyperplane because the Bernstein
evaluation matrix is invertible.  Thus the derivative of $\Phi_k$ is onto
$\mathcal V_k$, of dimension $(k-1)^2$.  Fixing complementary coordinates,
the inverse-function theorem provides a continuous local right inverse near
$0$ with fixed list size and interior coordinates.

For compact $K$, choose an even $L_0$ so large that $K/L_0$ lies in this
neighbourhood.  Realize $T/L_0$ locally and concatenate $L_0$ copies; by
additivity this realizes every $T\in K$ with one fixed even number $m$ of
points, continuously in $T$.  Choosing instead a sufficiently large odd
$L_1$ in the first construction gives
$Z_{\rm odd}=U_{L_1}\times U_{L_1}$, with odd size $h=L_1^2$ and
$\Phi_k(Z_{\rm odd})=0$.  All coordinates remain in one fixed interior
compact set.
\end{proof}

\subsection{The fixed-\texorpdfstring{$k$}{k} bounded-insertion theorem}
Fix $k\ge4$ and $\theta>0$.  A sequence is a
\emph{prescribed-aspect bounded-insertion completion sequence} if deleting
at most $M$ inserted points, for one constant $M$, leaves the two-sheet order
$\Pi(F)$ with
\[
 F=([a]\times[b])\cup(\{a+1\}\times R),\qquad
 a,b\to\infty,\quad a/b\to\theta,
\]
where $R\subseteq[b]$, $|R|=r$, and
\[
 \sup_{0\le z\le b}\left|\,|R\cap[z]|-\frac rbz\right|\le C
\]
for one constant $C$.  Insertion ranks are arbitrary, including endpoints;
no geometric separation or limiting location is assumed.

For $D\in\mathbb R^{k\times k}$ write
$\mathsf A_D(\tau)=\sum_{i=1}^kD_{i,\tau(i)}$, and define
\[
 q_k(\theta)=\min_{\substack{D\in\mathbb R^{k\times k}\\
                       \sum_{i,j}D_{ij}=0}}
       \max_{\tau\in\Ssym_k}|h_{k,\tau}(\theta)-\mathsf A_D(\tau)|.
\]
The minimum is attained because these assignment vectors form a
finite-dimensional closed subspace of $\mathbb R^{k!}$; it is the
$\ell_\infty$ distance from $(h_{k,\tau}(\theta))_\tau$.  Replacing $D$ by
$\mathcal P_kD\mathcal P_k$ preserves every assignment sum when
$\sum D_{ij}=0$.

\begin{theorem}[prescribed-aspect bounded-insertion optimum]\label{thm:fixed-k-insertion}
Fix $k\ge4$ and $\theta>0$. There is a constant $M_{k,\theta}$ and, for every sufficiently large
integer $N$, a permutation $\Xi_N^{(k,\theta)}\in\Ssym_N$ obtained from a
two-sheet row completion with $a/b\to\theta$ by at most $M_{k,\theta}$ point insertions such
that
\[
 \lim_{N\to\infty}\frac{\delta_{\Xi_N^{(k,\theta)},k}}{N^{k-1}}
 =\frac{q_k(\theta)}{2k!(k-2)!}.
\]
The family may depend on $k$ and $\theta$; the construction is nonconstructive. It
need not have rational or efficiently computable insertion coordinates.
Conversely, every prescribed-aspect bounded-insertion completion sequence has
\[
 \liminf_{\nu\to\infty}
 \frac{\delta_{\pi_\nu,k}}{N_\nu^{k-1}}
 \geq\frac{q_k(\theta)}{2k!(k-2)!}.
\]
This is an optimum only within the displayed completion class, not over all
permutations or all aspects.
\end{theorem}

\begin{proof}
Choose a total-sum-zero $D$ attaining $q_k(\theta)$.  For the upper
construction take the ceiling selector $R=R_{b,r}$ when $r>0$ and
$R=\varnothing$ when $r=0$, from Section~\ref{sec:construction}.  Its prefix
discrepancy is at most one, so
Lemma~\ref{lem:bounded-prefix-abel} and
Proposition~\ref{prop:boundary-leading} give a constant $C_k'$ such that,
for every resulting old order of even size $n=2|F|$, the matrix $E_n$ with
\[
 (E_n)_{k,t}=\frac{2}{(k-1)!}q_{k,t}(b,R),\qquad
 (E_n)_{i,t}=0\quad(i<k),
\]
has uniformly bounded entries.  Vandermonde gives
$\sum_{t=1}^kq_{k,t}(b,R)=0$, so $E_n$ has total sum zero and
$\mathsf A_{E_n}(\tau)$ is exactly the row term in
\eqref{eq:boundary-leading}.

Let $K$ be a compact ball in $\mathcal V_k$ containing all matrices
\[
 T_n=-s_k\mathcal P_k\left(\frac D{A_k}+E_n\right)\mathcal P_k.
\]
Lemma~\ref{lem:moment-realization} supplies one even list size $m$ realizing
all $T_n\in K$ and one odd neutral list of size $h$.  The ball is chosen
before $m,h$, so this choice is not circular.

For even $N$ use $n=N-m$; for odd $N$ use $n=N-m-h$.  Put $E=n/2$ and choose
\[
 b=\lfloor\sqrt{E/\theta}\rfloor,\qquad
 a=\left\lfloor\frac E b\right\rfloor,\qquad
 r=E-ab,
 \qquad R=R_{b,r}
\]
with $R=\varnothing$ if $r=0$.  For large sizes, $0\le r<b$ and
$a/b\to\theta$, while $|R\cap[z]|=\lfloor rz/b\rfloor$ gives a uniform
prefix bound.  Realize $T_n$ by $m$ insertions and append the $h$ neutral
insertions in the odd case.  Perturb the finitely many macro-coordinates by
$o(1)$ to make both total orders injective; choose any permutation for the
finitely many smaller sizes.

At the old size $n$, Proposition~\ref{prop:boundary-leading} gives, uniformly in $\tau$,
\[
 \frac{\#_\tau(\Pi(F))-\binom nk/k!}{n^{k-1}}
 =\frac{h_{k,\tau}(\theta)}{A_k}+\mathsf A_{E_n}(\tau)+o(1).
\]
By Lemma~\ref{lem:finite-padding}, the realization $\Phi_k(Z)=T_n$ changes
this vector by
$-\mathsf A_{D/A_k+E_n}(\tau)+o(1)$; centering does not change the assignment
sum of a total-sum-zero matrix.  The row term cancels, leaving
\[
 \frac{h_{k,\tau}(\theta)-\mathsf A_D(\tau)}{A_k}+o(1).
\]
Since $N=n+O_{k,\theta}(1)$, taking the maximum and then the limit proves
the upper equality in the theorem.

For the converse, apply Lemma~\ref{lem:finite-padding} to any sequence in
this class.  The bounded-prefix hypothesis controls the row term and all
errors uniformly, and the bounded insertion list contributes a centered
moment matrix of bounded mass.  After multiplying it by $A_k/s_k$ and
absorbing the row matrix $E_n$, the leading profile is
\[
 \frac{h_{k,\tau}(\theta)+\mathsf A_{D_\nu}(\tau)}{A_k}+o(1)
\]
for a total-sum-zero matrix $D_\nu$ (depending on $\nu$ and not required to
converge).  Replacing $D_\nu$ by $-D_\nu$ in the definition of $q_k(\theta)$
gives the pointwise lower bound $q_k(\theta)/A_k-o(1)$ on the maximum
discrepancy.  Taking the liminf proves the converse.
\end{proof}

\subsection{Exact LP constants and aspect optimization}
For each listed $(k,\theta)$ we give matching primal and dual certificates.  A
matrix $D$ is primal for $q$ if $\sum_{i,j}D_{ij}=0$ and, for every
$\tau\in\Ssym_k$,
\[
 \left|h_{k,\tau}(\theta)-\sum_iD_{i,\tau(i)}\right|\le q.
\]
A rational signed vector $(w_\tau)_{\tau\in\Ssym_k}$ is dual if
\[
 \sum_{\tau}|w_\tau|\le1,\qquad
 \sum_{\tau:\,\tau(i)=j}w_\tau
 \text{ is independent of }(i,j),
 \qquad
 \sum_{\tau}w_\tau h_{k,\tau}(\theta)=q.
\]
The constant-marginal condition annihilates assignment sums from every
total-sum-zero matrix, so these primal and dual conditions match by
$\ell_\infty$--$\ell_1$ duality.  An exact-arithmetic checker enumerates all
patterns for each listed $(k,\theta)$ over $\mathbb Q$ or $\mathbb Q(\sqrt3)$.
Its certificate arrays and checker are in the
\href{https://github.com/AutoclickerI/near-balanced-permutations/tree/v1.1.0}{validation repository},
file \texttt{validation/exact\_lp\_certificates.json}; from its root, run
\texttt{python3 -B -m validation.certificate\_audit}.  The six rows below
use the file's $(k,\theta)$ keys; the checker verifies these finite
certificates, while the all-size realization is the separate content of
Theorem~\ref{thm:fixed-k-insertion}.  Put
\begin{equation}\label{eq:k6-aspect-optimum}
 \rho=2-\frac{2\sqrt3}{3},\qquad q_*=\frac{111+5\sqrt3}{60}.
\end{equation}
\begingroup
\renewcommand{\arraystretch}{2}
\begin{equation}\label{eq:fixed-k-optima}
\begin{array}{c|c|c|c}
 k&\theta&q_k(\theta)&q_k(\theta)/(2k!(k-2)!)\\ \hline
 4&1&\dfrac{2}{3}&\dfrac{1}{144}\\
 5&1&\dfrac{5}{4}&\dfrac{1}{1152}\\
 6&1&2&\dfrac{1}{17280}\\
 7&1&\dfrac{62}{21}&\dfrac{31}{12700800}\\
 8&1&\dfrac{109}{28}&\dfrac{109}{1625702400}\\
 6&\rho&q_*&q_*/34560
\end{array}
\end{equation}
\endgroup
For a square-aspect dual, average $w_\tau$ with $w_{\tau^{-1}}$; its norm
is still at most one and its assignment marginals remain constant.  Since inversion swaps
$T_x,T_y$ and preserves $S$, this gives
\[
 q_k(\theta)\ge q_k(1)+\alpha_k(\theta+\theta^{-1}-2),\qquad
 \alpha_k=\frac12\sum_\tau w_\tau(T_x(\tau)+T_y(\tau)-2\mu_k).
\]
The supplied square duals have $\alpha_k=0,0,-1/5,13/42,3/8$ for
$k=4,\ldots,8$, respectively, so square aspect is optimal in these
completion classes for $k=4,5,7,8$.  For $k=6$, three supplied rational
duals give
\[
 q_6(\theta)\ge\max\left\{\frac85+\frac1{3\theta},\quad
               \frac85+\frac\theta3,\quad
               \frac{12}5-\frac{\theta+\theta^{-1}}5\right\}.
\]
Inversion also gives $q_k(\theta)=q_k(1/\theta)$.  On $0<\theta\le1$,
the first bound decreases and the third increases; they meet only at $\rho$,
where they equal $q_*$, so their maximum is strictly larger off $\rho$.
Since the exact primal at $\rho$ attains $q_*<2=q_6(1)$, the minimizing
aspects for $k=6$ are exactly $\rho$ and $1/\rho$ within the stated
completion framework.
For the four-pattern design below, write $B=B^{(4)}$ and
$K_\tau(u,v)=\frac1{36}\sum_{i=1}^4B_{i-1}(u)B_{\tau(i)-1}(v)-\frac1{144}$
for the single-insertion correction.

\subsection{Periodic row orientations}
\label{sec:oriented-four}
The optimum in Theorem~\ref{thm:fixed-k-insertion} fixes the standard
orientation of both sheets.  For four-pattern counts, changing the relative
orientation from row to row lowers the attainable constant.  Let
$\sigma=(\sigma_x)_{x\ge1}$ be a fixed periodic sequence in $\{-1,1\}$,
with period $p$ and mean $\eta=p^{-1}\sum_{x=1}^p\sigma_x$.  Replace the
second order in \eqref{eq:keys} by
\[
 K_Y^\sigma(x,y,c)=(y,\sigma_xcx,c),
 \qquad K_X(x,y,c)=(x,-cy,c).
\]
Write $\Pi^\sigma(F)$ for this order on positive-coordinate cell sets, and put
$U(\tau)=|\{i<4:|\tau(i+1)-\tau(i)|=1\}|$.  Thus $U$ counts pairs
adjacent in both position and value.  Define
\begin{equation}\label{eq:oriented-h}
 h^\eta_\tau(\theta)=\frac{T_x(\tau)-4/3}{\theta}
 +\theta(T_y(\tau)-4/3)+1-\frac{4\eta}{3}S(\tau)
 -\frac{2(1-\eta)}3U(\tau).
\end{equation}

\begin{lemma}[oriented leading profile]\label{lem:oriented-leading}
For fixed $\sigma$ and $a/b\to\theta>0$, the complete oriented rectangle
has centred four-profile $N^3h^\eta_\tau(\theta)/96+o(N^3)$, where $N=2ab$.
For its bounded-prefix partial-row completions and bounded insertions,
\eqref{eq:insertion-leading} holds at $k=4$ with $h_{4,\tau}$ replaced by
$h^\eta_\tau$.  All other terms and the uniformity in insertion ranks are
unchanged.
\end{lemma}
\begin{proof}
For compositions $\alpha,\beta\vDash4$, put $\ell=\ell(\alpha)$,
$a_i=p_\alpha(i)$ and $b_i=p_\beta(\tau(i))$.  Prescribe signs $\varepsilon_j$ for the distinct
rows.  Let $w^\varepsilon_\tau(\alpha,\beta)$ count colour words with
increasing keys $(-c_ib_i,c_i)$ within $X$-blocks and
$(\varepsilon_{a_i}c_ia_i,c_i)$ within $Y$-blocks, excluding repeated
signed points.  For each actual row choice $x_1<\cdots<x_\ell\le a$, set
$\varepsilon_j=\sigma_{x_j}$.  The exact core profile sums
$w^\varepsilon_\tau(\alpha,\beta)\binom b{\ell(\beta)}$ over these choices
and all compositions.  Defect-zero and defect-one weights are pattern-uniform for
every prescribed row-sign word: at a single pair the two constraints are
interchanged by reversing both colours, and singleton colours are free.
Their full contributions therefore vanish after centering, including all
lower coefficients of the coordinate-choice factors.

For defect two, average distinct-row signs independently with mean $\eta$.
The colour sums, aggregated over compositions with the indicated lengths,
are
\[
\begin{array}{c|c}
(\ell(\alpha),\ell(\beta))&\text{averaged colour sum}\\\hline
(2,4)&4(T_x+1)\\
(4,2)&4(T_y+1)\\
(3,3)&4\bigl(9-\eta S-(1-\eta)U/2\bigr).
\end{array}
\]
The one-axis rows follow from the triple and disjoint-pair counts in
Lemma~\ref{lem:defect-two}.  For the mixed row, use the active endpoints in that proof.
Edges with at most one common point have normalized weight $1$ if their
active endpoints differ, and $1-\eta h_e(i)v_f(i)$ otherwise.
If both endpoints coincide, the two points occupy one cell: the averaged
opposite-colour weight is $(1+\eta)/2$ for a descent and $(1-\eta)/2$ for
an ascent.  Summing these nine edge-pair cases gives
$9-\eta S-(1-\eta)U/2$.  The repository's
\texttt{validation/oriented\_four.py} also checks all $24$ patterns and
the polynomial coefficients in $\eta$ directly from the tie comparisons.

This independent-sign average gives the leading coefficient for a fixed
periodic sequence: for each residue word of $\ell$ increasing rows, the
number of choices is $a^\ell/(\ell!p^\ell)+O_p(a^{\ell-1})$.
Multiplying by $\binom b{\ell(\beta)}$ makes this error $O_p(N^{5/2})$
in defect two.  Larger defects have the same lower order.  For a uniform
four-pattern, $\mathbb E T_x=\mathbb E T_y=4/3$, $\mathbb ES=3/4$,
and $\mathbb EU=3/2$, the last identity counting the three adjacent pairs.
Center the table, divide by the coordinate factorials and by
$(2ab)^3$, and obtain \eqref{eq:oriented-h}/$96$.

For the partial row, replace the factor $\binom a{p-1}w_\tau$ in
\eqref{eq:partial-composition} by
\[
 \sum_{x_1<\cdots<x_{p-1}\le a}
 w^{(\sigma_{x_1},\ldots,\sigma_{x_{p-1}},\sigma_{a+1})}_\tau(\alpha,\beta).
\]
For $(p,q)=(4,4)$ the summand is $16$ for each row choice; for
$(3,4)$ or $(4,3)$ it is $8$, independently of the row signs and pattern.
Thus the clean Bernstein term and the defect-one density replacement in
Proposition~\ref{prop:boundary-leading} are unchanged.  Every remaining
term has size $O(a^{p-1}b^q)=O(b^5)=o(N^3)$ because $p+q\le6$.
For insertions, the leading coordinate-distinct term again ignores signs;
collision counts and old-rank displacement bounds depend only on row and
column sizes.  The proof of Lemma~\ref{lem:finite-padding} therefore gives
the asserted uniform expansion, including endpoints and moving ranks.
\end{proof}

\subsection{A cubic moment design}

Set $u_i=i/3$ for $0\le i\le3$, and define
\[
 D=\frac1{192}\begin{pmatrix}
 51&5&-8&-48\\
 18&-36&-63&81\\
 -63&-9&72&0\\
 -6&40&-1&-33
 \end{pmatrix},\qquad
 \beta_{is}=\frac34+\frac{D_{i+1,s+1}}2\quad(0\le s\le3).
\]
Each row of $D$ sums to zero.  Put
\[
 \mu_{i1}=\frac{\beta_{i1}+2\beta_{i2}+3\beta_{i3}}3,
 \qquad \mu_{i2}=\frac{\beta_{i2}}3+\beta_{i3},
 \qquad \mu_{i3}=\beta_{i3},
\]
and let $V_i$ consist of the three roots of
\[
 f_i(z)=z^3-\mu_{i1}z^2
       +\frac{\mu_{i1}^2-\mu_{i2}}2z
       -\frac{\mu_{i1}^3-3\mu_{i1}\mu_{i2}+2\mu_{i3}}6.
\]
Each $f_i$ changes sign from $-$ to $+$ on $(1/10,1/5)$,
from $+$ to $-$ on $(2/5,7/10)$, and from $-$ to $+$ on $(3/4,19/20)$.
These disjoint intervals contain all three roots, which are distinct and
lie in $(0,1)$.  Newton identities give
\[
 \sum_{v\in V_i}v^j=\mu_{ij}\quad(1\le j\le3),
 \qquad \sum_{v\in V_i}B_s(v)=\beta_{is}\quad(0\le s\le3).
\]
Use each pair $(u_i,v)$, $v\in V_i$, once, giving $12$ insertion
locations.  List them by increasing $i$ and then increasing root.

For odd orders append the five points in $H$, in lexicographic order, where
\[
 \begin{gathered}
 G_2=\left\{\frac12-\frac1{2\sqrt3},\ \frac12+\frac1{2\sqrt3}\right\},\\[-2pt]
 G_3=\left\{\frac12-\frac1{2\sqrt2},\ \frac12,\ \frac12+\frac1{2\sqrt2}\right\},\\[-2pt]
 H=(\{0\}\times G_2)\cup(\{1\}\times G_3).
 \end{gathered}
\]
The identities $\sum_{v\in G_2}B(v)=\frac12\mathbf1$ and
$\sum_{v\in G_3}B(v)=\frac34\mathbf1$ give $\mathcal P Q_H\mathcal P=0$.
The mass-$5$ identity then gives
$\sum_{(u,v)\in H}K_\tau(u,v)=\frac1{36}(\frac12+\frac34)-\frac5{144}=0$,
so this padding is neutral for every pattern.

Here is the exact moment calculation.  With $A_{s+1,i+1}=B_s(i/3)$ for $0\le s,i\le3$,
\[
 27A=\begin{pmatrix}27&8&1&0\\0&12&6&0\\0&6&12&0\\0&1&8&27\end{pmatrix},
 \qquad
 AD=\frac1{32}C,\qquad
 C=\begin{pmatrix}9&-1&-4&-4\\-1&-3&-2&6\\-4&-2&3&3\\-4&6&3&-5\end{pmatrix}.
\]
The last identity follows by multiplying the displayed integer matrices:
$(27A)(192D)=162C$.  Write
$\mathcal P=I_4-\mathbf1\mathbf1^{\mathsf T}/4$, where
$\mathbf1=(1,1,1,1)^{\mathsf T}$.  The $12$-point moment matrix is
$Q=A(\tfrac34\mathbf1\mathbf1^{\mathsf T}+D/2)$, so
$\mathcal P Q\mathcal P=C/64$.  Adding the neutral rule preserves this identity.  By
\eqref{eq:centered-assignment}, both lists have the total correction
\begin{equation}\label{eq:interior-correction}
 \sum_j K_\tau(u_j,v_j)=\frac1{2304}\sum_{i=1}^4C_{i,\tau(i)}.
\end{equation}

Choose the period-eight orientation
$(-,+,-,+,+,-,+,+)$, with mean $\eta=1/4$.
Use $m=12$ when $N$ is even and $m=17$ when $N$ is odd.
Set $\Sigma_N=\Pi_N$ if $N-m<100$.
Otherwise apply Lemma~\ref{lem:dimension-search} to $E=(N-m)/2$ with
$\lambda=3-\sqrt5$. Only its irrationality is needed for the profile limit.  If $r=0$, again use $\Pi_N$; this occurs
only finitely often.  In all other cases use the oriented order $\Pi^\sigma(F)$ with row
$R_{1/2}(b,r)$ and insert the listed $m$ points by
\eqref{eq:inserted-points}, keeping one global
index $j$ for all locations, including the five appended points at odd orders.
The final point count is exactly $N$.

\begin{theorem}[all-size four-point constant]\label{thm:sharp-four}
The explicit family $(\Sigma_N)$ satisfies
\[
 \lim_{N\to\infty}\frac{\delta_{\Sigma_N,4}}{N^3}=\frac1{192},
 \qquad
 \limsup_{N\to\infty}\frac{\delta_4(N)}{N^3}\le\frac1{192}.
\]
Conversely, let $\pi_\nu\in\Ssym_{N_\nu}$, $N_\nu\to\infty$, be row
completions with a fixed periodic orientation, a fixed limiting aspect
$\theta>0$, bounded prefix discrepancy and at most $M$ insertions,
where $M$ is fixed.  Then
\[
 \liminf_{\nu\to\infty}\frac{\delta_{\pi_\nu,4}}{N_\nu^3}\ge\frac1{192}.
\]
This optimum allows all fixed periods and positive limiting aspects, but
is not asserted over arbitrary permutations.
\end{theorem}

\begin{proof}
The dimension search gives $a/b\to1$ and $r/b\to3-\sqrt5$.
By Lemmas~\ref{lem:midpoint-quadrature} and~\ref{lem:oriented-leading}, the
row term vanishes; \eqref{eq:interior-correction} gives, for both parities,
\[
 L_\tau=\frac{h^{1/4}_\tau(1)}{96}
       +\frac1{2304}\sum_{i=1}^4 C_{i,\tau(i)}.
\]
Substitution in the local formula gives the following complete ledger:
\[
\begin{array}{c|l}
2304L_\tau&\tau\\\hline
-12&1243,1324,2143,3412,4231\\
-8&4321\\
-4&2341,3421,4123,4312\\
0&1234,1432,2314,2431,3124,3214,4132\\
12&1342,1423,2134,2413,3142,3241,4213
\end{array}
\]
The limiting maximum is $12/2304=1/192$.

For optimality, define the ten-pattern functional
\[
\begin{split}
16\mathcal W(z)={}&-2z_{1243}+2z_{1342}+2z_{1423}+2z_{2134}-z_{2143}\\
 &+z_{2413}+z_{3142}+2z_{3241}-z_{3412}+2z_{4213}.
\end{split}
\]
Its coefficient $\ell_1$ norm is $1$, and every position--value pair has
coefficient sum $1/8$.  Hence it annihilates assignment sums of total-sum-zero
matrices, including the centred row and insertion corrections.  Directly,
\[
 \mathcal W(1)=\tfrac12,\quad
 \mathcal W(T_x)=\mathcal W(T_y)=\tfrac34,\quad
 \mathcal W(S)=\tfrac18,\quad\mathcal W(U)=\tfrac14.
\]
Thus \eqref{eq:oriented-h} gives
\[
 \mathcal W((h^\eta_\tau(\theta))_\tau)
 =\frac13+\frac{\theta+\theta^{-1}}{12}\ge\frac12,
\]
independently of $\eta$.  Lemma~\ref{lem:oriented-leading} and
$\|\mathcal W\|_1=1$ give the lower bound $1/(2\cdot96)=1/192$.
The ledger attains this bound.
\end{proof}


\medskip\noindent\textbf{Concluding remarks.}
The exponent-optimal family is $k$-independent, while completion constants
depend on the allowed orientations. The relative-balancing range of this
family is sharp. Global constants, optimal aspects beyond the tabulated orders,
and growing-pattern ranges attainable by other families remain open.

\medskip\noindent\textbf{Declarations.} This work received no funding. The authors
declare no financial or non-financial conflicts of interest.

\noindent\textbf{Data availability.} No external datasets are used.
The validation source is at
\url{https://github.com/AutoclickerI/near-balanced-permutations/tree/v1.1.0}.
It includes exact primal/dual certificates over $\mathbb Q$ and
$\mathbb Q(\sqrt3)$ and finite regression checkers; the analytic proofs
are in the article. From the repository root run \texttt{bash validation/run.sh}.

\noindent\textbf{Use of generative AI.} OpenAI's GPT-5.6 and GPT-6 Astra
assisted literature searches, proof exploration and checking, auxiliary
computations, and LaTeX and English editing.  The authors
take responsibility for the mathematical arguments, sources, computations, and
text.

\begin{thebibliography}{99}
\small
\setlength{\emergencystretch}{2em}
\setlength{\itemsep}{0pt}

\bibitem{BLL}
G.~Beniamini, N.~Lavee, and N.~Linial,
\newblock How balanced can permutations be?,
\newblock \emph{Combinatorica} 45 (2025), Article 9,
\newblock \href{https://doi.org/10.1007/s00493-024-00127-x}
 {\texttt{doi:10.1007/s00493-024-00127-x}}.

\bibitem{Cooper}
J.~N. Cooper,
\newblock Quasirandom permutations,
\newblock \emph{J. Combin. Theory Ser. A} 106 (2004), 123--143,
\newblock \href{https://doi.org/10.1016/j.jcta.2004.01.006}
 {\texttt{doi:10.1016/j.jcta.2004.01.006}}.

\bibitem{CooperArithmetic}
J.~N. Cooper,
\newblock Quasirandom arithmetic permutations,
\newblock \emph{J. Number Theory} 114 (2005), 153--169,
\newblock \href{https://doi.org/10.1016/j.jnt.2005.05.003}
 {\texttt{doi:10.1016/j.jnt.2005.05.003}}.

\bibitem{CooperDukesFuzzy}
J.~Cooper and P.~J. Dukes,
\newblock Fuzzy latin squares and balanced permutation pattern statistics,
\newblock preprint (2026),
\newblock \href{https://arxiv.org/abs/2608.05335}
 {\texttt{arXiv:2608.05335}}.

\bibitem{CooperPetrarca}
J.~N. Cooper and A.~Petrarca,
\newblock Symmetric and asymptotically symmetric permutations,
\newblock preprint (2008),
\newblock \href{https://arxiv.org/abs/0801.4181}
 {\texttt{arXiv:0801.4181}}.

\bibitem{HoppenEtAl}
C.~Hoppen, Y.~Kohayakawa, C.~G. Moreira, B.~R\'ath, and
R.~M. Sampaio,
\newblock Limits of permutation sequences,
\newblock \emph{J. Combin. Theory Ser. B} 103 (2013), 93--113,
\newblock \href{https://doi.org/10.1016/j.jctb.2012.09.003}
 {\texttt{doi:10.1016/j.jctb.2012.09.003}}.

\bibitem{KralPikhurko}
D.~Kr\'al' and O.~Pikhurko,
\newblock Quasirandom permutations are characterized by 4-point densities,
\newblock \emph{Geom. Funct. Anal.} 23 (2013), 570--579,
\newblock \href{https://doi.org/10.1007/s00039-013-0216-9}
 {\texttt{doi:10.1007/s00039-013-0216-9}}.

\bibitem{SeymourZaslavsky}
P.~D. Seymour and T.~Zaslavsky,
\newblock Averaging sets: A generalization of mean values and spherical designs,
\newblock \emph{Advances in Mathematics} 52 (1984), 213--240,
\newblock \href{https://doi.org/10.1016/0001-8708(84)90022-7}
 {\texttt{doi:10.1016/0001-8708(84)90022-7}}.

\end{thebibliography}

\end{document}
